\documentclass[11pt,a4paper]{article}
\usepackage[margin=27mm]{geometry}
\usepackage[T1]{fontenc}
\usepackage[utf8]{inputenc}
\usepackage{lmodern,amsmath,amssymb,amsthm,mathtools,booktabs,longtable,array}
\usepackage[dvipsnames]{xcolor}
\usepackage{microtype}
\usepackage[hidelinks]{hyperref}
\usepackage{enumitem}
\setlist{itemsep=3pt,topsep=5pt}
\setlength{\parskip}{4pt}
\setlength{\parindent}{0pt}
\setlength{\emergencystretch}{3em}
\newtheorem{theorem}{Theorem}[section]
\newtheorem{lemma}[theorem]{Lemma}
\newtheorem{proposition}[theorem]{Proposition}
\newtheorem{corollary}[theorem]{Corollary}
\newtheorem{definition}[theorem]{Definition}
\DeclareMathOperator{\CH}{CH}
\DeclareMathOperator{\Pic}{Pic}
\DeclareMathOperator{\NS}{NS}
\DeclareMathOperator{\End}{End}
\DeclareMathOperator{\Hom}{Hom}
\DeclareMathOperator{\Alb}{Alb}
\title{JAC1964: primitive incidence polarizations,\newline pencil heights, and the arithmetic complement}
\author{Run JAC1964--20261121\\GPT-6 Astra Pro; collaborative mathematical investigation}
\date{21 November 2026\\Working manuscript: checkpoint version}
\begin{document}
\maketitle
\begin{abstract}
Accepting the supplied 032 results, we prove that the correctly signed
primitive incidence self-duality of an odd-dimensional abelian variety is
a rational polarization in every characteristic. An explicit quadric
pencil refutes the proposed uncorrected H9 height identity, even when
all fibres are geometrically integral and the restricted cycles are
rationally zero. We give the admissible local correction and a precise
no-correction criterion. Algebraic Lefschetz correspondences prove the
restriction and fixed-part assertions of the sixth theorem. An
all-primes monodromy argument also determines the variable incidence
sign under explicit Lefschetz monodromy hypotheses. The arithmetic
investigation constructs an actual orthogonal Chow Picard summand for
CM varieties and isolates its complementary height problem. A Fermat
quartic Ceresa class shows that this complement is nonempty.
The unrestricted seventh theorem, with the customary meaning of
$\tau$-equivalence, is false even for CM abelian threefolds.
This is a checkpoint manuscript; the statement-level audit is in
\texttt{CHECK\_EVENTS.tsv}. Its arithmetic additions are still being
extended and reviewed. No general arithmetic Hodge-index theorem is
claimed.
\end{abstract}
\tableofcontents
\newpage
\section{Scope, inputs, and conventions}
All new deductions are labelled by stable finding IDs. ``Proved'' means
proved from the explicitly accepted and cited inputs, with the
whole-proof reading recorded in the check log; it does not mean formal
verification. The model effort setting is not exposed in this session.
The original archive images were not supplied, so the transcription
audit is mathematical, not palaeographic.

The incidence interpretation of the higher Picard variety is specified
below. We work with rational cycles and rational homomorphisms of
abelian varieties, clearing denominators on actual families. We do not
tensor geometric point groups with $\mathbf Q$ to define varieties.
No integral universal cycle or principal polarization is asserted.

The attached 032, Theorem 1.1 and Corollary 8.4, the supplied
van Hoften--Sempliner and Milne reports, and Alp\"oge's advertised
Riemann-hypothesis and Dirichlet-GRH statements are accepted as inputs,
as requested. Their full proofs and the omitted Lean dependency tree
were not independently certified here. The bridge sketch and H9 are
questions to test, not input theorems. The exact scopes of the analytic
inputs are recorded in the arithmetic section.

\part{Checked findings and cited background}
% Proof-ready contributions by /root/programme_audit.
% Run JAC1964-20261121; GPT-6 Astra Pro; effort setting not exposed.
% All statements are rational, unless explicitly stated otherwise.

\section{Incidence Picard varieties and the next theorem of the programme}

For a smooth projective variety $V$ of dimension $d$ over an algebraically
closed field, write $A^p(V)=\mathrm{CH}^p(V)_{\mathrm{alg}}$ and
\[
 I^p(V)=\{z\in A^p(V):\Gamma_*z=0\text{ in }\mathrm{Pic}^0(T)
 \text{ for every smooth projective }T
 \text{ and }\Gamma\in\mathrm{CH}^{d-p+1}(V\times T)\}.
\]
Saito's higher Picard variety, denoted $P^p(V)$ here, represents
$A^p(V)/I^p(V)$ on algebraically closed points
\cite[Theorem~(3.6)(i)]{Saito1979}.
This is the incidence construction suggested by cote~44, scans~2 and~9.
It must not be silently substituted for the universal regular algebraic
representative, a different universal problem.  Saito supplies duality
\emph{up to isogeny}, not an assertion of a canonical integral principal
polarization \cite[Proposition~(4.4)]{Saito1979}.
Rational families below mean rational homomorphisms of abelian varieties,
obtained by clearing the denominators of the cycles first.  They do not
mean that one tensors the abstract group of geometric points with
$\mathbf Q$; that operation would lose the construction over
$\overline{\mathbf F}_p$, where all abelian-variety points are torsion.

\begin{lemma}[JAC-P00: rational descent of the incidence construction]
\label{jac:P00}
Let $F$ be any field and let $V/F$ be smooth and projective.  The
incidence Picard variety of $V_{\overline F}$, together with its quotient
map on algebraically trivial cycles, has a model over $F$ in the isogeny
category.  If $T/F$ is abelian and $u:T\to P^p(V)$ is a rational
homomorphism over $F$, there is a cycle
$Z\in\mathrm{CH}^p(T\times V)_{\mathbf Q}$ whose family map is $u$.
\end{lemma}

\begin{proof}
Saito's quotient is canonical: its universal property among Picard
homomorphisms gives a unique isomorphism from each conjugate quotient
to the quotient of $V_{\overline F}$.  These isomorphisms satisfy the
cocycle identity.  The variety, its group law, the finite correspondence
data witnessing that its quotient map is a Picard homomorphism, and
one surjective generating family are defined over a finite extension
of $F$.  The latter can be taken from Saito's Poincar\'e cycle, whose
family map is a positive integer times the identity.  An automorphism
fixing these data fixes the quotient map on that generating family.
Uniqueness of the comparison isomorphism, and surjectivity of the
family map, then force that comparison to be the identity.  This
proves, using finite algebraic data rather than the set of all cycles,
that the canonical descent datum is trivial on an open subgroup of
the absolute Galois group.

Here are the details that also cover an imperfect $F$.  Enlarge a
finite separable extension to a finite Galois extension $L/F$, and
choose $r\ge0$ so that the preceding finite field of definition lies
in $L^{1/p^r}$ in characteristic $p$; take $r=0$ in characteristic
zero.  Raising the coefficients of the variety and its group law to
$p^r$ gives an abelian variety $Q/L$, and relative Frobenius gives an
isogeny $f:P_{\overline F}\to Q_{\overline F}$.
Transport the canonical descent datum by $f$ and its rational inverse.
For $\sigma\in\operatorname{Gal}(L/F)$ this gives
\[
 c_\sigma:{}^\sigma Q\longrightarrow Q
 \quad\text{in the rational category},\qquad
 c_{\sigma\tau}=c_\sigma\,{}^\sigma c_\tau.
\]
These maps are defined over $L$: homomorphisms of abelian varieties
acquire no new elements under purely inseparable extensions, and the
remaining Galois invariance follows from the cocycle identity and the
triviality of the datum on automorphisms fixing $L$.  For the former
standard rigidity fact, see \cite[Example~4.2.14, p.~38]{ConradAV}.

Rational Galois descent is effective by the following explicit
construction.  Let $G=\operatorname{Gal}(L/F)$ and
$R=\operatorname{Res}_{L/F}Q$.  Over $L$ identify
$R_L=\prod_{\sigma\in G}{}^\sigma Q$.  The maps
\[
 i:Q\longrightarrow R_L,\quad
 i(x)=(c_\sigma^{-1}x)_\sigma,
 \qquad
 q:R_L\longrightarrow Q,\quad
 q((x_\sigma))=\frac1{|G|}\sum_\sigma c_\sigma x_\sigma
\]
satisfy $qi=1$.  The cocycle identity makes $e=iq$ Galois invariant,
so $e\in\operatorname{End}^0_F(R)$.  For a positive integer $N$ with
$Ne$ integral, its image is an abelian variety over $F$ whose base
change is isomorphic to $P$ in the isogeny category.  This supplies the
required rational model and its transported quotient map.  Uniqueness
in the rational category follows from the same quotient universal
property.

Finally pull a normalized rational Poincar\'e cycle back by $u$ over
$\overline F$.  The resulting family is defined over some finite
extension $F'/F$.  Put $d_s=[F':F]_{\rm sep}$ and
$d_i=[F':F]_{\rm insep}$, and let
$\pi:(T\times V)_{F'}\to T\times V$.
Define $Z=(d_sd_i)^{-1}\pi_*Z_{F'}$.  After base change to
$\overline F$, flat base change for Chow groups expresses this as
the average of the $d_s$ conjugate cycles, each with multiplicity
$d_i$.  The multiplicity is the length of the corresponding local
factor of $F'\otimes_F\overline F$.  Every conjugate induces $u$,
because $u$ is defined over $F$ for the compatible quotient descent
datum.  Hence $Z$ induces $u$.  This explains explicitly both the
separable conjugates and the inseparable factor; it does not treat a
purely inseparable extension as if it had extra automorphisms.
\end{proof}

\begin{lemma}[JAC-P01: detection by the odd realization]
\label{jac:P01}
Let $V,W$ be smooth projective varieties over an algebraically closed
field, and let
$\Gamma\in\mathrm{CH}^{\dim V-p+q}(V\times W)_{\mathbf Q}$.
Fix a prime $\ell$ invertible in the field.  If the action of $\Gamma$
from $H^{2p-1}(V,\mathbf Q_\ell)$ to
$H^{2q-1}(W,\mathbf Q_\ell(q-p))$ is zero, then the induced rational
homomorphism $P^p(V)\longrightarrow P^q(W)$ is zero.

More generally, for an abelian variety $T$ and
$Z\in\mathrm{CH}^q(T\times W)_{\mathbf Q}$, the homomorphism
$T\longrightarrow P^q(W)$ defined by $t\mapsto Z_t-Z_0$ depends only
on the K\"unneth component
\[
 [Z]_{1,2q-1}\in H^1(T,\mathbf Q_\ell)\otimes
 H^{2q-1}(W,\mathbf Q_\ell(q)).
\]
In particular, it is zero if that component is zero.
\end{lemma}

\begin{proof}
For the second assertion, test $Z_t-Z_0$ by an arbitrary incidence
correspondence
$\Xi\in\mathrm{CH}^{\dim W-q+1}(W\times S)$.
The composition $\Xi\circ Z$ is a divisor class on $T\times S$.
Its component in $H^1(T)\otimes H^1(S)$ depends only on
$[Z]_{1,2q-1}$, by the K\"unneth formula and integration over $W$.
A divisor correspondence on $T\times S$ gives a homomorphism
$T\to\mathrm{Pic}^0(S)$; its rational Tate-module action is precisely
the mixed degree-one component.  Faithfulness of the rational Tate
module on homomorphisms of abelian varieties therefore says that this
homomorphism is zero when the indicated mixed component is zero.
Consequently every $Z_t-Z_0$ is annihilated by every incidence tester.
It is zero in $A^q(W)/I^q(W)$ by definition, proving the assertion.
The argument applies to rational cycles after clearing denominators.

Every algebraically trivial cycle is a finite sum of differences of
fibres of families parametrized by smooth projective curves.
For one such family
$U\in\mathrm{CH}^{p}(C\times V)$, its only relevant odd component
gives the map $H^1(C)\to H^{2p-1}(V)(p-1)$.
Composition with $\Gamma$ is zero by the hypothesis.  Composing further
with each incidence tester on $W$ gives a divisor correspondence on
$C\times S$ with zero mixed component.  The induced
$\operatorname{Alb}(C)\to\mathrm{Pic}^0(S)$ is zero, by the same
faithfulness argument.  Thus $\Gamma_*$ sends every element of
$A^p(V)$ into $I^q(W)$, and its action on $P^p(V)$ is zero.
\end{proof}

The lemma does not assume the standard conjecture~D.  The use of
cohomology concerns correspondences acting on families of cycles;
it does not identify homological and incidence equivalence of individual
cycles.

\begin{theorem}[JAC-P02: Theorem~6 for incidence Picard varieties]
\label{jac:P02}
Let $X$ be smooth projective of dimension $n$ over an algebraically
closed field, and assume that its normalized Lefschetz lowering
operator is algebraic.  For a smooth hyperplane section $j:Y\hookrightarrow X$,
restriction induces
\[
 j^*:P^p(X)\longrightarrow P^p(Y).
\]
It is an isogeny if $2p\le n-1$, has finite kernel if $2p=n$, and
is surjective if $2p\ge n+1$.  Each assertion means the corresponding
split statement in the rational category of abelian varieties, hence
the asserted statement for the original integral homomorphism.

Suppose further that $Y=Y_\eta$ is the generic member of a pencil with
smooth codimension-two base locus, and put $K=k(t)$.  In the critical
case $n=2p$, the image of $P^p(X)_K$ is the image of the $K/k$-trace
of $P^p(Y_\eta)$, up to isogeny.  Thus the rational form of all clauses
of cote~44, scan~4, Theorem~6, holds whenever the Lefschetz standard
conjecture holds for $X$.  In particular, it holds for abelian $X$
in every characteristic.
\end{theorem}

\begin{proof}
Write $L$ for multiplication by the hyperplane class and $\Lambda$
for the normalized lowering operator: it sends $L^a v$ to
$L^{a-1}v$ for a primitive vector $v$ and $a>0$, and kills $v$.
The standard conjecture supplies an algebraic correspondence for this
operator.  The correspondence
\[
 R=\Lambda_X\circ j_*
\]
has degree zero and induces $P^p(Y)\to P^p(X)$.
Put $r=2p-1$.
For $r<n$, no Lefschetz string terminates in degree $r$, so
\[
 Rj^*=\Lambda_X L_X=\operatorname{id}\quad\text{on }H^r(X).
\]
By Lemma~\ref{jac:P01}, the same identity holds on $P^p(X)$.
If $r<n-1$, weak Lefschetz makes $j^*:H^r(X)\to H^r(Y)$ an
isomorphism, hence $j^*R$ is also the identity on $H^r(Y)$ and on
$P^p(Y)$.  This proves the first assertion and, when $r=n-1$,
the assertion of finite kernel.

For $r\ge n$, the Gysin map
$j_*:H^r(Y)\to H^{r+2}(X)(1)$ is an isomorphism by the dual of
weak Lefschetz.  No Lefschetz string begins in degree $r+2>n$,
so $L_X\Lambda_X=\operatorname{id}$ there.  For $y\in H^r(Y)$,
\[
 j_*j^*Ry=L_X\Lambda_Xj_*y=j_*y.
\]
Injectivity of $j_*$ gives $j^*R=\operatorname{id}$ on $H^r(Y)$;
Lemma~\ref{jac:P01} transfers this identity to $P^p(Y)$.
These are precisely the algebraic inverse mechanisms of cote~16,
scan~33, formula~B(4), and scans~40 and~45, clauses~(iv) and~(iv bis).

Here is a cycle-level proof of the fixed-part assertion, so no inference
from invariant cohomology to algebraicity of a correspondence is left
unstated.  Write $p=m$, $n=2m$, let $B$ be the base locus, and let
\[
 b:\widetilde X=\operatorname{Bl}_B X\longrightarrow X,
 \qquad e:E\hookrightarrow\widetilde X,
 \qquad q:E\longrightarrow B
\]
be the blow-up, exceptional divisor, and its projection.  Let $T/k$
be an abelian variety and
$u:T_K\longrightarrow P^m(Y_\eta)$ a homomorphism.  Saito's
Poincar\'e $m$-cycle supplies, after division by its positive integer
multiplier, a rational cycle on $P^m(Y_\eta)\times Y_\eta$
inducing the identity \cite[\S4.1]{Saito1979}.  Lemma~\ref{jac:P00}
descends its pullback by $u$ to a rational cycle
\[
 Z_\eta\in\mathrm{CH}^m(T_K\times Y_\eta)_{\mathbf Q}
\]
inducing $u$.  In particular the construction is valid even when
$K=k(t)$ is imperfect.  Subtracting $Z_0$ normalizes
the family at the identity of $T$.

Close a cycle representative in $T\times\widetilde X$.
The codimension-two blow-up formula gives
\[
 \overline Z=(1_T\times b)^*Z_X+
 (1_T\times e)_*(1_T\times q)^*Z_B,
\]
where
$Z_X\in\mathrm{CH}^m(T\times X)_{\mathbf Q}$ and
$Z_B\in\mathrm{CH}^{m-1}(T\times B)_{\mathbf Q}$ are defined over
$k$.  Write $k_B:B\hookrightarrow X$ and $i:B_K\hookrightarrow Y_\eta$.
The mixed component relevant to $Z_B$ is
$H^1(T)\otimes H^{n-3}(B)$.
Weak Lefschetz identifies $H^{n-3}(X)$ with $H^{n-3}(B)$.
An algebraic inverse is
\[
 S=\Lambda_X^2\circ(k_B)_*.
\]
Indeed, $(k_B)_*k_B^*=L_X^2$ and $\Lambda_X^2L_X^2$ is the
identity on degree $n-3$.  Define an algebraic rational cycle over $k$
by
\[
 W_X=Z_X+L_X\circ S\circ Z_B\in
 \mathrm{CH}^m(T\times X)_{\mathbf Q}.
\]
For $x\in H^{n-3}(X)$ the projection formula gives
\[
 i_*k_B^*x=i_*i^*j_\eta^*x=L_{Y_\eta}j_\eta^*x
 =j_\eta^*L_Xx.
\]
It follows that $Z_\eta$ and $j_\eta^*W_X$ have the same
$H^1(T)\otimes H^{n-1}(Y_\eta)$ component.  Lemma~\ref{jac:P01}
says they induce the same rational homomorphism from $T_K$ to
$P^m(Y_\eta)$.  The latter factors through $P^m(X)_K$ since $W_X$
is defined over $k$.  When $m=1$ the exceptional term has no mixed
component, since $H^{-1}(B)=0$, and the same conclusion follows.

Every homomorphism from a constant abelian variety has therefore image
in $\operatorname{Im}j_\eta^*$.  Conversely, this image is itself the
image of a homomorphism from the constant variety $P^m(X)_K$.
The universal property of the $K/k$-trace proves the asserted equality
of images up to isogeny, using the rational $K$-model of
Lemma~\ref{jac:P00}.
\end{proof}

\section{An explicit test of the primitive sign}

\begin{proposition}[JAC-P03: the uncorrected correspondence is indefinite]
\label{jac:P03}
There exist normalized algebraic families of codimension-two cycles
on an abelian threefold for which the divisor correspondence of
cote~44, scan~9, is nondegenerate and indefinite.  Neither quotienting
by its kernel nor multiplication by a common sign makes it a
polarization.  The primitive qualification in scan~8 is essential.
\end{proposition}

\begin{proof}
Let $E$ be an elliptic curve, set $X=E_1\times E_2\times E_3$,
write $a_i$ for the pullback of $[0]\in\mathrm{CH}^1(E_i)$, and use
$h=a_1+a_2+a_3$ as polarization.  On $E\times E$ put
\[
 P=[\Delta_E]-[E\times0]-[0\times E].
\]
As a correspondence $P$ is the symmetric degree-one Chow--K\"unneth
projector, hence ${}^tP\circ P=P$.
As a divisor correspondence it induces
\[
 E\longrightarrow\widehat E,
 \qquad t\longmapsto\mathcal O_E(t-0)=-\lambda_{\mathcal O_E(0)}(t),
\]
where the polarization convention is
$\lambda_L(t)=t_t^*L\otimes L^{-1}$.
Thus the sign calibration is fixed, rather than hidden in a convention.

Define two families on $E_t\times X$ by
\[
 z_\pm=P_{t,1}\,(a_2\pm a_3).
\]
Their incidence divisor correspondences are
\[
 D_\pm=\deg_{E_2\times E_3}(a_2\pm a_3)^2\,P
       =\pm2P,
\]
and the mixed incidence correspondence between the two families
is zero because $(a_2+a_3)(a_2-a_3)=0$.
Use independent parameters for the two families, obtaining a family
over $T=E\times E$.  Its symmetric homomorphism $T\to\widehat T$
has diagonal blocks $-2\lambda_{\mathcal O(0)}$ and
$+2\lambda_{\mathcal O(0)}$.  It is an isogeny, but its associated
rational line bundle restricts to a negative-degree line bundle on one
elliptic factor and a positive-degree one on the other.  It cannot be
ample.  Its negative cannot be ample either.  Passing through a finite
quotient does not change this conclusion, since the pullback of an
ample line bundle by an isogeny is ample.

The degree-one parameter components make the primitive issue explicit.
For $\alpha\in H^1(E_1)$, they have images
\[
 u_\pm(\alpha)=\alpha(a_2\pm a_3).
\]
One has $u_+(\alpha)=L_X\alpha$ and $L_Xu_-(\alpha)=0$.
Moreover, for $\alpha,\beta\in H^1(E_1)$,
\[
 \int_Xu_+(\alpha)u_+(\beta)=2\int_E\alpha\beta,
 \qquad
 \int_Xu_-(\alpha)u_-(\beta)=-2\int_E\alpha\beta,
\]
and the cross pairing vanishes.  The lower Lefschetz and primitive
pieces require different signs.  In the divisor convention above,
the primitive $m=2$ correspondence $D_-=-2P$ already gives the
positive homomorphism $2\lambda_{\mathcal O(0)}$.
\end{proof}

\section{Primitive middle cycles need not reach the incidence Jacobian}

\begin{proposition}[JAC-P04: an abelian-fourfold domain test for H9]
\label{jac:P04}
Let $X=E_1\times E_2\times E_3\times E_4$ over $\mathbf C$,
$a_i=[0]$ pulled back from $E_i$, and $h=a_1+a_2+a_3+a_4$.
For a sufficiently high-degree Lefschetz pencil relative to $h$, let
$Y_\eta$ be its geometric generic member.  The codimension-two cycle
\[
 \alpha=(a_1-a_2)(a_3-a_4)
\]
satisfies $h\alpha=0$ and $\deg\alpha^2=4$.  Its restriction to
$Y_\eta$ is homologically trivial, but no nonzero multiple of that
restriction is algebraically equivalent to zero.
\end{proposition}

\begin{proof}
The relations $a_i^2=0$ give
\[
 h\alpha=0,
 \qquad
 \alpha^2=(-2a_1a_2)(-2a_3a_4)=4a_1a_2a_3a_4.
\]
In particular $\alpha$ has nonzero cohomology class.  Since the Gysin
map $H^4(Y_\eta)\to H^6(X)$ is an isomorphism by weak Lefschetz
and duality, the equality $j_*j^*\alpha=d\,h\alpha=0$ gives
$j^*\alpha=0$ in cohomology, for the degree-$d$ pencil.

Apply SGA~7~II, Expos\'e~XX, Theorem~4.1
\cite[printed p.~355]{SGA7II}.  Its hard-Lefschetz hypothesis holds
for $X$.  Its algebraic-retraction hypothesis~B is supplied by the
algebraic Lefschetz operator, as noted in Expos\'e~XX, \S3.4.
For sufficiently high degree the pencil satisfies condition~A by
Expos\'e~XVIII, Corollary~6.4.  In characteristic zero its variable
middle cohomology has level greater than one for sufficiently high
degree, by Expos\'e~XX, Remark~4.3 (printed p.~356).
The cited theorem says that a middle-dimensional cycle on $X$ with
algebraically trivial generic restriction has zero cohomology class
on $X$.  Applied to any nonzero multiple of $\alpha$, this contradicts
$\deg\alpha^2=4$.  Therefore no such multiple is algebraically
trivial on $Y_\eta$.
\end{proof}

This example identifies the enlargement between cote~44, scan~4,
Theorem~5, which explicitly requires algebraic triviality on $Y$,
and H9, which assumes only homological triviality.  The latter condition
does not provide a point of the incidence Picard variety.  It provides
a class in the Griffiths group in this example.  A general intermediate
Jacobian or normal function is a different object and does not remove
this distinction.

\begin{proposition}[JAC-P05: the customary reading of Theorem~7 is false]
\label{jac:P05}
If $\tau$-equivalence means that a nonzero integral multiple is
algebraically equivalent to zero, the unrestricted equality of numerical
and $\tau$-equivalence in cote~44, scan~6, fails already for complex
abelian threefolds of CM type, including one defined over a number field.
\end{proposition}

\begin{proof}
Take the Fermat quartic
\[
 C:\quad T_0^4+T_1^4=T_2^4
\]
and embed it in its Jacobian $J=J(C)$ using a rational base point.
Bloch's Theorem~(4.1), printed p.~104 of \cite{Bloch1984L}, states
that its Ceresa cycle $\Xi=[C]-[-1]_*[C]$ has infinite order modulo
algebraic equivalence.  This is the rational non-torsion assertion,
stronger than merely $\Xi\not\sim_{\rm alg}0$.
Bloch's Lemma~(3.1), printed p.~99, gives an isogeny
\[
 J\longrightarrow E\times E\times E',\qquad
 E:y^2=x^3-x,\quad E':y^2=x^3+x.
\]
Both elliptic curves have CM by $\mathbf Q(i)$, as is seen from
$(x,y)\mapsto(-x,iy)$; thus $J$ is of CM type.
On $H^4(J,\mathbf Q)$ the involution $[-1]^*$
acts as $(-1)^4=1$.  Thus this cycle is homologically trivial, hence
numerically trivial by compatibility of cycle classes and intersection.
Its infinite order modulo algebraic equivalence means that it is not
$\tau$-equivalent to zero.  This gives the required counterexample.
\end{proof}

Ceresa's original theorem also gives a very general genus-three
Jacobian example.  The explicit Bloch example has the advantage that
both the infinite-order statement and the CM isogeny have directly
verified primary-source locators.  A theorem making Hodge classes
algebraic says nothing that would force a homologically trivial
cycle such as $\Xi$ to become algebraically trivial.

For divisors the statement is correct in every characteristic:
SGA~6, Expos\'e~XIII, Theorem~4.6 identifies numerical triviality
with $\tau$-triviality; definitions are in \S\S4.1 and~4.4.
The supplied archive page does not itself define $\tau$.  Accordingly,
the preceding refutation is explicitly tied to this customary reading,
not to a guessed replacement of $\tau$ by homological equivalence.

\begin{lemma}[JAC-P06: the finite-field meaning of the last claim]
\label{jac:P06}
Over $k=\overline{\mathbf F}_p$, every algebraically trivial Chow class
is torsion.  Hence $\tau$-triviality means torsion in the Chow group,
and numerical equivalence equal to $\tau$-equivalence would identify
the rational Chow group with cycles modulo numerical equivalence.
\end{lemma}

\begin{proof}
An algebraically trivial cycle is a finite sum of images, under cycle
correspondences, of degree-zero zero-cycles on smooth projective curves.
For a curve $C/k$, its degree-zero Chow group is $J(C)(k)$.  Every
point and the curve are defined over a finite field, and each finite
group $J(C)(\mathbf F_q)$ is torsion.  Thus $J(C)(k)$ is torsion,
and so is each of its images under a correspondence.  A finite sum of
torsion elements is torsion.  If a nonzero multiple of a Chow class
is algebraically trivial, another nonzero multiple is therefore zero
in the Chow group.  The converse is immediate.
\end{proof}


\section{The primitive incidence polarization in every characteristic}
\label{sec:T4}

All varieties in this section are over an algebraically closed field $k$.
We work with rational correspondences and with $\ell$-adic cohomology,
$\ell\ne\operatorname{char}k$; Tate twists in displayed K\"unneth components
are suppressed when no confusion is possible. A rational polarization means
an element of $\operatorname{NS}(B)_{\mathbf Q}$ a positive integral multiple
of which is the class of an ample line bundle. It does not mean a principal
polarization or an integral universal cycle.

Write $P^p(X)=\operatorname{Pic}^p(X)$ for Saito's higher Picard variety:
\[
 A^p(X)/K^p(X)\simeq P^p(X)(k),\qquad
 A^p(X)=\operatorname{CH}^p(X)_{\rm alg}.
\]
Here $K^p$ is the incidence kernel, tested by divisor-valued
correspondences. The construction, rational Poincar\'e cycles and the
canonical duality isogeny are the cited results
\cite[Theorem~3.6, \S4.1, Proposition~4.4]{Saito1979}.
The notation is the precise incidence interpretation of cote~44,
scan~2, Theorems~1--3, and scan~8. It must not be replaced without proof
by a universal regular algebraic representative.

\begin{lemma}[JAC-T4.1: primitive incidence summand]
\label{lem:primitive-summand}
Let $A$ be an abelian variety of dimension $n=2m-1$, and let $h$ be a
rational ample divisor class. Let $e_h$ denote the cohomological projector
of $H^n(A)$ onto
\[
 H^n(A)_{\rm prim}=\ker\bigl(L_h:H^n(A)\longrightarrow H^{n+2}(A)\bigr)
\]
along $L_hH^{n-2}(A)$, extended by zero in other degrees.
It is represented by a rational algebraic correspondence. It induces an
idempotent on $P^m(A)$ in the isogeny category. Its image $B=P^m(A)_{\rm prim}$
has a nondegenerate symmetric rational incidence correspondence
$d_B\in\operatorname{NS}(B\times B)_{\mathbf Q}$, rigidified along the two
axes. For a rational Poincar\'e family $z\in\operatorname{CH}^m(B\times A)_{\mathbf Q}$
inducing the inclusion $B\longrightarrow P^m(A)$, it is
\begin{equation}\label{eq:incidence}
 d_B=\operatorname{pr}_{12*}(z_{13}\cdot z_{23})_{11}.
\end{equation}
The subscript $11$ removes the pullbacks of divisor classes on the two
factors. The family may be chosen with cohomology entirely in
$H^1(B)\otimes H^n(A)_{\rm prim}$.
\end{lemma}

\begin{proof}
The Lefschetz standard conjecture and the algebraic K\"unneth projectors
hold for abelian varieties; for an exact classical locator see
\cite[Theorem~2A11]{Kleiman68}. Thus the primitive projector is algebraic.
The incidence-detection lemma Lemma~\ref{jac:P01} says that a correspondence
whose relevant odd cohomological action is zero acts as zero on higher
Picard varieties. In particular, its action depends only on the
homological class of that component. Applying this to $e_h^2-e_h$ gives
an actual idempotent in $\operatorname{End}^0(P^m(A))$.

The Poincar\'e pairing on $H^n(A)$ makes its primitive summand orthogonal
to $L_hH^{n-2}(A)$: if $L_hu=0$, then
$\int_Au\,L_hv=\int_A(L_hu)\,v=0$.
Hard Lefschetz makes this an orthogonal direct sum of nondegenerate
spaces. Consequently $e_h$ is self-transpose in degree $n$. Saito's
transpose compatibility for incidence duality
\cite[Proposition~4.4, especially (4.4.1)--(4.4.2)]{Saito1979}, together
with incidence detection, says that the induced idempotent is
self-adjoint for the canonical rational duality of $P^m(A)$.
Writing the resulting isogeny decomposition as
$P^m(A)\sim B\times B'$, this duality is block diagonal. Since the full
duality is an isogeny, each block is an isogeny. In particular the block
on $B$ is nondegenerate.

Choose Saito's Poincar\'e cycle and divide it by its positive integer
$k_m$ so that it induces the identity, then pull it back to $B\times A$.
Apply the degree-one K\"unneth projector on the parameter factor and
$e_h$ on the $A$ factor. These operations preserve the required map to
the incidence quotient and give the asserted mixed cohomological
component. Subtracting constant restrictions rigidifies the family at
the origin without changing this component. Composition of
correspondences gives \eqref{eq:incidence}. Equivalently, the
unnormalized universal cycle would multiply this correspondence by
$k_m^2$, a positive rational square. Symmetry follows by interchanging
the two parameter factors. The block argument proves its
nondegeneracy. This also proves independence of the chosen normalized
family, in the N\'eron--Severi group, from Saito's incidence identities.
\end{proof}

\begin{lemma}[JAC-T4.2: transfer of the curve test to an abelian product]
\label{lem:curve-test}
Let $C$ be a smooth projective pointed curve of positive genus $b$, and
let $A$ have dimension $2m-1$. Suppose a rational algebraic class
$\xi\in H^1(C)\otimes H^{2m-1}(A)$ satisfies
$(1\otimes L_h)\xi=0$. Let $J=J(C)$, let
$a:C\longrightarrow J$ be its Abel map, and let $\theta$ be its canonical
principal polarization. There is a rational algebraic class
$w\in H^1(J)\otimes H^{2m-1}(A)$ with
$(a\times1)^*w=\xi$. For $H=\theta+h$ on $J\times A$,
\begin{equation}\label{eq:curve-transfer}
 H^bw=0,\qquad
 \int_{J\times A}w^2H^{b-1}
 =(b-1)!\int_{C\times A}\xi^2.
\end{equation}
Hence 032, Corollary~8.4, gives
\[
 (-1)^m\int_{C\times A}\xi^2\ge0,
\]
with equality if and only if $\xi=0$.
\end{lemma}

\begin{proof}
The Abel map identifies $h^1(C)$ and $h^1(J)$ by rational algebraic
correspondences. One can use its graph, the divisor correspondence of
the Poincar\'e bundle, and the inverse Lefschetz correspondence on $J$
to obtain an algebraic inverse to $a^*:H^1(J)\to H^1(C)$.
More explicitly, $w=(b-1)!\,(L_\theta^{b-1})^{-1}(a\times1)_*\xi$ on the indicated mixed component; the inverse is algebraic by the Lefschetz standard conjecture for $J$. Indeed, $a_*a^*x=x\theta^{b-1}/(b-1)!$ for $x\in H^1(J)$. Apply this inverse to the first factor of $\xi$ and project to
$H^1(J)\otimes H^{2m-1}(A)$; the latter projector is algebraic because
both varieties in the product are abelian. This constructs $w$ as an
algebraic class, not merely as a cohomological tensor. It is unique as a
cohomological tensor, and $h\,w=0$.

The dimension of $J\times A$ is $b+2m-1$, and $w$ has codimension $m$.
Every term in $(\theta+h)^bw$ with a positive power of $h$ vanishes.
The remaining term $\theta^bw$ also vanishes because its degree on $J$
is $2b+1$. Thus $w$ is primitive for the product ample class.
Likewise $w^2H^{b-1}=w^2\theta^{b-1}$.
The canonical polarization of the Jacobian gives
\[
 a_*[C]=\frac{\theta^{b-1}}{(b-1)!}
 \quad\hbox{in cohomology}.
\]
Only its pairing with $H^2(J)$ is needed here. In this form it follows
immediately from the identification of the canonical alternating
pairings on $H^1(C)$ and $H^1(J)$: in a symplectic basis
$\theta=\sum_{i=1}^b x_i\wedge y_i$, expansion of
$\theta^{b-1}$ gives $(b-1)!$ times the complementary wedges.
The identification of the pairings is Jacobian autoduality
\cite[Chapter~III, Theorem~6.8 and Summary~6.11]{MilneAV}.
Projection formula now gives \eqref{eq:curve-transfer}.

Apply the Hodge standard conjecture in the exact form of 032,
Corollary~8.4, to the abelian variety $J\times A$, codimension $m$ and
ample class $H$. Its degree exponent is
$(b+2m-1)-2m=b-1$, and its sign is $(-1)^m$.
The factor $(b-1)!$ is positive. The equality assertion follows from
positive definiteness and the injectivity of $a^*$ in degree one.
\end{proof}

\begin{theorem}[JAC-T4.3: cote 44, Theorem 4, on the primitive incidence quotient]
\label{thm:T4}
Accept 032, Corollary~8.4, as in the task specification. For every
abelian variety $A$ of odd dimension $2m-1$ over an algebraically closed
field of any characteristic, and every rational ample class $h$, the
symmetric homomorphism on $P^m(A)_{\rm prim}$ defined by
\[
 (-1)^m d_B
\]
is a rational polarization. Equivalently, with
$\phi_L(x)=t_x^*L\otimes L^{-1}$, the rational divisor class
\begin{equation}\label{eq:T4-divisor}
 \ell_B=\frac{(-1)^m}{2}\,\Delta_B^*d_B
 \quad\in\operatorname{NS}(B)_{\mathbf Q}
\end{equation}
is ample. After multiplication by a positive integer it defines an
actual polarization in every characteristic, including characteristic
$2$.
\end{theorem}

\begin{proof}
Use Lemma~\ref{lem:primitive-summand} to choose a normalized universal
family $z$ and the nondegenerate symmetric incidence correspondence
$d_B$. For a symmetric rigidified divisorial correspondence, the
class $\Delta^*d_B/2$ is the rational N\'eron--Severi class inducing
that correspondence. Indeed, for a line bundle $L$ the associated
biextension class is
$m_B^*c_1(L)-p_1^*c_1(L)-p_2^*c_1(L)$, whose diagonal pullback is
$[2]^*c_1(L)-2c_1(L)=2c_1(L)$.
Thus \eqref{eq:T4-divisor} induces $(-1)^md_B$, which is still
nondegenerate.

Let $C\to B$ be the normalization of any integral curve in $B$.
It is nonconstant, so $C$ has positive genus. Translate it so that a
chosen point maps to the origin; translations do not change
N\'eron--Severi classes. Pulling back $z$ gives $z_C$ on $C\times A$,
whose cohomology is in
$H^1(C)\otimes H^{2m-1}(A)_{\rm prim}$.
Refined pullback and projection formula for the incidence product give
\begin{equation}\label{eq:degree-test}
 \deg_C(\ell_B)=\frac{(-1)^m}{2}
       \deg_{C\times A}(z_C\cdot z_C).
\end{equation}
There are no omitted degree-zero or degree-two parameter components:
the family was projected to the degree-one component before this
calculation. Lemma~\ref{lem:curve-test} makes the right side
nonnegative. Hence $\ell_B$ is nef.

A nef line bundle on an abelian variety whose associated symmetric
homomorphism is an isogeny is ample. To recall a short precise version
of this criterion, clear denominators and call the resulting bundle
$L$. Nondegeneracy gives $\chi(L)\ne0$ and
$\deg\phi_L=\chi(L)^2$, whereas nefness and Riemann--Roch give
$\chi(L)>0$. For an ample $H_0$, the polynomial
$\chi(L+tH_0)=(c_1(L)+tc_1(H_0))^{\dim B}/(\dim B)!$
is positive for all $t\ge0$. Mumford's index theorem therefore gives
index zero; its effective nondegenerate criterion then gives
ampleness \cite[Chapter~III, \S16, Index Theorem, and Chapter~II,
\S6, Application~1]{MumfordAV}. This argument is valid in arbitrary
characteristic. Applying it to \eqref{eq:T4-divisor} proves the theorem.
The zero-dimensional $B$ case is immediate. Rational coefficients
justify all divisions even when the characteristic divides a chosen
integral denominator; no assertion about an integral principal
polarization follows.
\end{proof}

\paragraph{Archive meaning and the exact input.}
This is the primitive reading required by the marginal instruction in
cote~44, scan~8, and the incidence product in scan~9. Cote~16,
scans~34, 36, 38, 40 supplies the Lefschetz language; scans~49--50 places
positivity in the ring of correspondences. The passage from cycle
positivity to an ample divisor is the curve test above; it is not the
unproved assertion that a cohomologically nondegenerate correspondence
is automatically a polarization. The input 032 is used only on
$J(C)\times A$ in Lemma~\ref{lem:curve-test}. No Hodge conjecture for a
non-CM characteristic-zero abelian variety and no new algebraicity
claim for such Hodge classes is inferred.

\paragraph{Sign calibration.}
For an elliptic curve, the normalized divisor
$P=\Delta-E\times\{0\}-\{0\}\times E$ induces
$x\mapsto\mathcal O_E(x-0)=-\phi_{\mathcal O_E(0)}(x)$.
Its self-incidence is $P$, and $\deg\Delta^*P=-2$.
Thus the sign in \eqref{eq:T4-divisor} is indeed $(-1)^m$ under the
stated convention: for $m=1$ it gives the positive elliptic
polarization. The frequently written sign $(-1)^{m-1}$ belongs instead
to the alternating primitive cohomological pairing before the
incidence-to-divisor identification. Confusing these two conventions
would reverse the conclusion already for a curve.


% Research contribution of /root/h9_height, JAC1964-20261121.
% This is a section fragment, not a standalone document.
% Sources are identified in H9_sources.md. All claims below have an explicit
% logical status. No formal proof certificate is asserted.

\section{The pencil formula: a counterexample and its local repair}

The text of cote 44, scan 4, assumes that the restrictions of the two
cycles are algebraically equivalent to zero. H9 replaces this hypothesis by
homological triviality. These are different hypotheses. Moreover, geometric
integrality of the fibres does not remove higher-codimension vertical
corrections. The following counterexample satisfies even the stronger
condition that the restrictions are rationally equivalent to zero.

\subsection*{H9-C1. A quadric counterexample, with both local contributions}

\textbf{Claim.} For every integer $m\geq2$, over every algebraically closed
field $k$, there is a
Lefschetz pencil on a smooth $2m$-dimensional quadric $X$, all of whose fibres
are geometrically integral, and a codimension-$m$ cycle $Z$ such that
\[
 Z|_{Y}=0\quad\text{in }\operatorname{CH}^{m}(Y),\qquad
 (-1)^m\deg(Z^2)=2.
\]
Here $Y$ is the generic fibre over $K=k(t)$. Consequently every
bilinear pairing on its Chow group gives height zero to the restriction,
whereas the uncorrected expression in H9(a) gives $2$. The claim also
refutes the literal formula of cote 44, scan 4, under just its recorded
geometrically integral-fibre hypothesis.

\textbf{Proof.}
Write
\[
 X=\{uv+\textstyle\sum_{i=1}^{m} w_i z_i=0\}
       \subset\mathbf P^{2m+1}.
\]
Use the pencil of hyperplanes $a u+b v=0$, with parameter $[a:b]$.
Its base is the smooth quadric
\[
 B=\{u=v=0,\ \textstyle\sum w_i z_i=0\}\simeq Q^{2m-2}.
\]
The members with $ab\ne0$ are smooth odd-dimensional quadrics. This remains
true in characteristic two: after eliminating $v$, the equation is
$-t u^2+\sum w_i z_i=0$; the only possible projective singular point
from its partial derivatives would have all $w_i,z_i$ zero, and then
the equation forces $u=0$ since $t\ne0$. The two
remaining members, $u=0$ and $v=0$, are cones over $B$, each with precisely
one ordinary quadratic singularity, at the $v$- or $u$-vertex respectively.
For $m\geq2$ the base $B$ is a smooth integral positive-dimensional quadric,
so its cones are integral. The axis is transverse, and the total space
obtained by blowing up $B$ is smooth. At each singular point the local
equation is $\sum w_i z_i=0$ in $2m$ variables; its polar form is
nondegenerate also in characteristic two. Thus this is a Lefschetz pencil in
the precise sense of SGA 7 II, XVII, 2.2. It lies in the generic orbit of
pencils under the orthogonal group: equivalently its pencil line meets the
dual quadric transversely in two points.

Put
\[
 L_+=\mathbf P\langle w_1,\ldots,w_m\rangle,\qquad
 L_-=\mathbf P\langle w_1,\ldots,w_{m-1},z_m\rangle
       \subset B.
\]
These are maximal $(m-1)$-planes in the two different families on $B$.
Let $P_{\infty,+},P_{\infty,-}$ be their cones with the $u$-vertex and
$P_{0,+},P_{0,-}$ their cones with the $v$-vertex. They are $m$-planes on
$X$. Denote their Chow classes by the same symbols and set
\[
 Z=P_{\infty,+}-P_{\infty,-}.
\]
The two planes in this difference belong to the two different families on
$X$. Replacing the vertex while retaining $L_+$ interchanges these families;
the same holds for $L_-$. This assertion can be checked without a
classification theorem. Quotient by
$\langle w_1,\ldots,w_{m-1}\rangle$ and use the rank-one matrix description
of the quadric surface $uv+w_mz_m=0$, with matrix
$\left(\begin{smallmatrix}u&w_m\\-z_m&v\end{smallmatrix}\right)$.
The two row lines are fibres of one ruling and the two column lines are
fibres of the other. Varying the row (respectively column) gives a
$\mathbf P^1$-family joining $P_{\infty,+}$ to $P_{0,-}$ (respectively
$P_{\infty,-}$ to $P_{0,+}$), after adjoining the common subspace.
Consequently
\[
 [P_{0,+}]-[P_{0,-}]=-Z.
\]

We check the intersection number directly. For an $m$-plane $P\subset X$,
the normal-bundle sequence is
\[
 0\longrightarrow N_{P/X}\longrightarrow
 \mathcal O_{\mathbf P^m}(1)^{\oplus(m+1)}
 \longrightarrow\mathcal O_{\mathbf P^m}(2)\longrightarrow0.
\]
The last map is surjective since $X$ is smooth. Therefore
\[
 [P]^2=[h^m]\frac{(1+h)^{m+1}}{1+2h}
       =\frac{1+(-1)^m}{2}.
\]
Here $[h^m]$ denotes extraction of the coefficient of $h^m$. If $\alpha$
and $\beta$ are the two middle plane classes, a linear section degenerating
to two planes gives $h^m=\alpha+\beta$. Since $\deg(h^m\alpha)=1$,
\[
 \alpha\beta=\frac{1-(-1)^m}{2},\qquad
 (\alpha-\beta)^2=2(-1)^m.
\]
For the degeneration $h^m=\alpha+\beta$, intersect $X$ with the
$\mathbf P^{m+1}$ on which all $z_i$ vanish. Its equation is $uv=0$, and
its components are $P_{\infty,+}$ and $P_{0,+}$; the row--column families
above identify $[P_{0,+}]=\beta$. Thus every Chow relation used in this
calculation has an explicit geometric realization.

On the generic fibre, take the affine parameterization $v=-t u$. Its
equation is $-t u^2+\sum w_i z_i=0$. The restrictions of the two planes
appearing in $Z$ are $L_+$ and $L_-$. They are rationally equivalent over
$K$, not merely after passing to $\overline K$. Indeed, put
$W=\langle w_1,\ldots,w_{m-1}\rangle$ and consider the smooth $K$-conic
\[
 C_t=\{-t u^2+w_mz_m=0\}\subset\mathbf P^2_K.
\]
It is parameterized by
\[
 [r:s]\longmapsto[u:w_m:z_m]=[rs:r^2:t s^2].
\]
The family of projective $(m-1)$-planes
$\mathbf P(W\oplus\ell)$, for the isotropic lines $\ell$ represented by
$C_t$, lies in $Y$. At $[r:s]=[1:0]$ and $[0:1]$ its members are $L_+$
and $L_-$. Since these are $K$-rational points of $C_t\simeq\mathbf P^1_K$,
the two fibre cycles are rationally equivalent. Thus $Z|_Y=0$ as asserted.

We now compute the correction at each singular fibre. Write
$q:\widetilde X=\operatorname{Bl}_{B}X\to X$,
$j:E\hookrightarrow\widetilde X$, and $p:E\to B$. Set
\[
 t_X=q^*Z,\quad \delta_B=[L_+]-[L_-],\quad
 e=j_*p^*\delta_B,
\]
and let $u_\infty,u_0$ be the differences of the strict transforms of
$P_{\infty,\pm}$ and $P_{0,\pm}$, respectively. Then
\[
 u_\infty=t_X-e,\qquad u_0=-t_X-e.                 \tag{H9.1}
\]
For completeness, these identities follow directly from the incidence
model $\widetilde X\subset X\times\mathbf P^1$. On
$P_{\infty,\pm}\times\mathbf P^1$, the incidence equation is $a u=0$;
its divisor is the union of the plane in the fibre $a=0$ and
$L_\pm\times\mathbf P^1$. This is precisely the refined pullback of
$P_{\infty,\pm}$. On $P_{0,\pm}\times\mathbf P^1$ the equation is
$b v=0$, yielding the analogous identity. Taking differences proves
(H9.1), including the coefficients and signs.

The projection formula gives $t_X e=0$ because $q_*e=0$. Since
$\mathcal O_E(E)=\mathcal O_E(-1)$ and $p_*c_1(\mathcal O_E(1))=1$,
\[
 e^2=-\delta_B^2=-2(-1)^{m-1}=2(-1)^m,
 \qquad t_X^2=2(-1)^m.
\]
It follows that
\[
 \begin{array}{c|cc}
  &u_\infty&u_0\\\hline
 u_\infty&4(-1)^m&0\\
 u_0&0&4(-1)^m
 \end{array},\qquad
 t_Xu_\infty=2(-1)^m,\quad t_Xu_0=-2(-1)^m.
                                                        \tag{H9.2}
\]
In particular $t_X=(u_\infty-u_0)/2$ is itself a rational vertical cycle.
Its admissible correction is
\[
 t_X-\tfrac12u_\infty+\tfrac12u_0=0.
\]
Each local Schur-complement term is
$(2(-1)^m)^2/(4(-1)^m)=(-1)^m$. Thus each singular fibre contributes $1$
after multiplication by the H9 sign $(-1)^m$, and the two contributions
cancel the uncorrected value $2$. This proves the counterexample and the
claimed local refinement.\hfill$\square$

This is a mathematical insufficiency of the stated hypothesis, not evidence
that the word ``int\`egres'' was transcribed incorrectly. Without images of
the original manuscript no palaeographic alteration is justified.

\subsection*{H9-T1. The corrected intersection formula}

Let $f:\mathcal X\to C$ be a projective morphism from a smooth projective
$2m$-fold to a smooth projective curve over an algebraically closed field,
with smooth generic fibre $Y$. Use rational Chow groups. Call
$x\in\operatorname{CH}^m(\mathcal X)_{\mathbf Q}$ admissible if
$x|_Y$ is numerically trivial and
\[
 \deg(xv)=0
 \quad\text{for every codimension-$m$ cycle $v$ supported in fibres.}
                                                        \tag{H9.3}
\]
The geometric intersection-height normalization used here is
\[
 h_m(x|_Y,y|_Y)=(-1)^m\deg(xy)                    \tag{H9.4}
\]
for admissible extensions. Its independence of extensions follows from
the localization exact sequence: two extensions of the same Chow class
differ, modulo rational equivalence, by a vertical cycle, and (H9.3)
annihilates this difference. This is the signed degree of the pairing in
Kahn, \emph{Refined height pairing}, Definition 2.2, Proposition 2.8 and
equation (2.9). The terminology ``geometric intersection height'' specifies
the construction; comparison with a particular Bloch biextension or a
N\'eron--Tate pairing is a separate compatibility statement.

Suppose $x_0,y_0$ extend numerically trivial generic-fibre cycles and that
vertical corrections are to be made inside a finite-dimensional space
$W$ spanned by numerical classes of vertical cycles. Let
$R=\operatorname{rad}(W)$ for the intersection form. Assume $x_0$ and $y_0$
are orthogonal to $R$. Choose representatives $v_1,\ldots,v_r$ for a basis
of $W/R$, put $M_{ij}=\deg(v_iv_j)$ and
\[
 a_i=\deg(x_0v_i),\qquad b_i=\deg(y_0v_i).
\]
Then $M$ is invertible, the corrected classes are
\[
 x=x_0-\sum_i(M^{-1}a)_i v_i,\qquad
 y=y_0-\sum_i(M^{-1}b)_i v_i,
\]
and direct bilinear expansion gives
\[
 h_m(x|_Y,y|_Y)=(-1)^m\left(\deg(x_0y_0)-a^{\!t}M^{-1}b\right).
                                                        \tag{H9.5}
\]
The expression is independent of the chosen representatives modulo $R$.
If $W$ contains all the vertical numerical classes, the corrected cycles
satisfy (H9.3). Conversely, existence of a correction in $W$ requires
$x_0\perp R$, because adding an element of $W$ cannot change its pairing
with $R$. The finite-dimensional linear algebra proves this criterion is
also sufficient. Hence (H9.5) states both the formula and its exact
existence gate.

For a pencil with blowup map $q$ and $x_0=q^*Z$, $y_0=q^*Z'$, the
projection formula gives $\deg(x_0y_0)=\deg(ZZ')$. Fibres at different
closed points are disjoint, so the matrix can be organized into local
blocks, after separating the radical. Formula (H9.5) is the requested
local variant, including higher-dimensional vertical cycles; it is not
restricted to irreducible components of fibres. For $m=1$ it recovers the
usual intersection-matrix correction in Shioda, \emph{On the Mordell--Weil
lattices}, Theorem 8.6, equations (8.11)--(8.15), printed p.~228. The
original printed page was rendered and read in this run.

\subsection*{H9-T2. A cohomological hypothesis that removes all corrections}

\textbf{Claim.} Let a Lefschetz pencil on a smooth projective $2m$-fold
$X$ have generic fibre $Y$ and smooth blowup $q:\widetilde X\to X$.
Fix $\ell\ne\operatorname{char}k$. Suppose at each singular fibre $Y_s$
the specialization map
\[
 H^{2m}(Y_s,\mathbf Q_\ell(m))\longrightarrow
 H^{2m}(Y_{\bar\eta_s},\mathbf Q_\ell(m))       \tag{H9.6}
\]
is injective. If $Z,Z'$ restrict to homologically trivial cycles on $Y$,
then $q^*Z,q^*Z'$ are admissible and
\[
 h_m(Z|_Y,Z'|_Y)=(-1)^m\deg(ZZ').               \tag{H9.7}
\]
This statement holds in every characteristic subject to the explicit
injectivity hypothesis (H9.6).

\textbf{Proof.} On the smooth locus of the pencil the fibre classes of
$q^*Z$ form a section of the locally constant sheaf
$R^{2m}f_*\mathbf Q_\ell(m)$. The section is zero at the geometric generic
point, hence zero on this locus. At a singular fibre, compatibility of the
class of a global cycle with specialization and (H9.6) imply
$i_s^*\operatorname{cl}(q^*Z)=0$ in the ordinary cohomology of $Y_s$,
where $i_s:Y_s\hookrightarrow\widetilde X$. For any $m$-dimensional cycle $V$
inside any fibre, the compatibility of the cycle class with cap products
and the trace gives
\[
 \deg(q^*Z\cdot V)=
 \operatorname{Tr}_{V}(i_s^*\operatorname{cl}(q^*Z)\cap[V])=0.
\]
Thus $q^*Z$ is admissible; similarly for $Z'$. Formula (H9.7) follows
from (H9.4) and the projection formula.\hfill$\square$

In the ordinary nodal Picard--Lefschetz setting, a sufficient condition
for (H9.6) is that every local vanishing cycle is nonzero in the middle
cohomology of the nearby smooth fibre. Indeed, the relevant exact sequence
is
\[
 H^{2m-1}(Y_{\bar\eta_s})\longrightarrow\mathbf Q_\ell(-m)
 \longrightarrow H^{2m}(Y_s)
 \longrightarrow H^{2m}(Y_{\bar\eta_s})\longrightarrow0.
\]
The first arrow is pairing with the vanishing cycle. Poincar\'e duality
makes it surjective when that cycle is nonzero, proving injectivity of the
next specialization map. This identifies the additional local condition
missing in H9; the ordinary quadric pencil has zero middle cohomology and
therefore zero vanishing cycles. The precise SGA framework is SGA 7 II,
XVIII, 5.3.2 (printed p.~287), 5.3.4--5.3.5 (p.~288): condition $(A)$ says
that the direct-image sheaves agree with the extensions of their
restrictions to the smooth locus. Proposition 6.2 (p.~319) identifies
$(A)$ with nontrivial inertia action on middle cohomology at every
singular fibre. These original printed pages were independently read in
the programme audit. Using condition $(A)$, rather than an unqualified
quadratic-monodromy convention, makes (H9.6) available in every
characteristic covered by that condition.

Assuming the attached 032 Corollary 8.4, if $X$ is an abelian $2m$-fold,
the right side of (H9.7) is positive definite on its primitive algebraic
cohomology. If $h$ is the ample hyperplane class defining the pencil, put
\[
 C_{\rm prim}=\{Z\in\operatorname{CH}^m(X)_{\mathbf Q}:
                          h\operatorname{cl}(Z)=0\},\quad
 R=\operatorname{im}(C_{\rm prim}\to\operatorname{CH}^m(Y)_{\mathbf Q}),
 \quad
 R_0=\operatorname{im}(\operatorname{CH}^m_{\rm hom}(X)_{\mathbf Q}
                         \to\operatorname{CH}^m(Y)_{\mathbf Q}).
\]
Here $\ker(j^*:H^{2m}(X)\to H^{2m}(Y))=\ker L_h$: indeed
$j_*j^*=L_h$, and $j_*:H^{2m}(Y)\to H^{2m+2}(X)$ is injective by
Poincar\'e duality and weak Lefschetz in degree $2m-2$. Thus every element
of $C_{\rm prim}$ satisfies the restriction hypothesis of H9-T2.
For pencils satisfying (H9.6), $h_m$ on $R$ has radical exactly $R_0$:
the inclusion of $R_0$ in the radical follows because a homologically
trivial ambient class has zero intersection with every algebraic class;
the opposite inclusion follows from (H9.7) and positive definiteness of
the primitive form. In particular the restriction map induces an
isometry
\[
 C_{\rm prim}/\operatorname{CH}^m_{\rm hom}(X)_{\mathbf Q}
       \xrightarrow{\sim}R/R_0
\]
for the signed primitive intersection form on the left and the geometric
height on the right. For injectivity, if a restriction belongs to $R_0$,
subtract a homologically trivial ambient lift; the resulting restriction
is zero, so its height is zero and positivity forces its ambient
cohomology class to vanish. This proves the assertion without pretending
that a cohomology class has a canonical Chow representative. The argument
does not turn homologically trivial restrictions into points of an
algebraic intermediate Jacobian.

\subsection*{H9-T3. The precise quotient statement behind H9(b)}

Here is the valid linear-algebra implication, separating all geometric
inputs. Let $V$ be a rational vector space of cycle classes equipped with
a symmetric form $B$, and let
\[
 a:V\longrightarrow
 (J(K)/\operatorname{Tr}_{K/k}J(k))\otimes\mathbf Q
\]
be a defined homomorphism. Suppose $J$ is polarized and a height comparison
has established
\[
 B(v,w)=\langle a(v),a(w)\rangle_{\rm NT}.
                                                        \tag{H9.8}
\]
Then $B$ is positive semidefinite and
\[
 \operatorname{rad}B=\ker a.
\]
Indeed the right side is positive definite on the Lang--N\'eron space
after quotienting torsion, so $B(v,v)=0$ is equivalent to $a(v)=0$; the
bilinear identity proves the radical statement. Thus $B$ is positive
definite on $V$ if and only if $a$ is injective; it is always positive
definite on $V/\ker a$.

Applied to a pencil, (H9.8) must use the corrected form (H9.5). Replacing
it by the uncorrected primitive intersection form additionally requires
that every local correction vanish. The quadric example has nonzero
primitive classes with zero generic restriction, so injectivity on the
original primitive space fails. This does not by itself refute a bare
logical equivalence of two positivity assertions (both are true for a
quadric); it refutes the claimed identification and its asserted proof.

There is also a distinct domain obstruction. H9 assumes only homological
triviality, whereas Grothendieck's $P^m$ is constructed from algebraically
trivial cycles. SGA 7 II, XX, Theorem 4.1 and Remark 4.3 provide examples
where these restrictions are not algebraically trivial, even with
$X=E^4$: with $a_i$ the divisor class from the $i$th elliptic factor and
$h=\sum a_i$, the class
\[
 \alpha=(a_1-a_2)(a_3-a_4)
\]
satisfies $h\alpha=0$ and $\alpha^2=4$. On sufficiently high-degree
Lefschetz hyperplane sections, its restriction is homologically trivial
but not algebraically trivial. The original locators are printed
pp.~355--356; the full source hypotheses and their verification are
supplied by the independent programme audit. In particular, the regular
homomorphism from algebraically trivial cycles defining $P^m$ cannot be
applied to this cycle. One can ask for other, additional maps or
projections on a larger group, but H9 has not specified or proved such a
construction. Thus its displayed ``image'' in the algebraic $P^m$ of
cote 44 is not defined on the stated domain by the construction cited.

The repaired route must retain three pieces: the admissible height on
homologically trivial cycles; the local vertical form; and, only on the
algebraically representable part, comparison with an abelian
N\'eron--Tate height. The supplied CM-Hodge and standard-conjecture theorems
give positivity of the primitive intersection form for abelian $X$; they
do not remove either of the two explicit counterexamples above.

\subsection*{H9-T4. Polarization of the fixed primitive part of an ample section}

\textbf{Claim.} Assume 032, Corollary 8.4. Let $X$ be an abelian variety
of dimension $2m$ over an algebraically closed field of any characteristic,
let $j:Y\hookrightarrow X$ be a smooth ample divisor, and write
$h=[Y]$, $n=2m-1$. The fixed primitive summand of the incidence higher
Picard variety $P^m(Y)$ has a rational polarization given by its
$(-1)^m$-signed incidence correspondence. Its cohomological subspace is
\[
 F=j^*\ker\bigl(L_h^2:H^{n}(X)\longrightarrow H^{n+4}(X)\bigr)
       \subset H^n(Y).
\]
This proves the relevant part of Theorem 4 even though $Y$ need not be
abelian. The residual summand is the variable part, on which $j_*$ is zero.

\textbf{Proof.}
Hard Lefschetz on $X$ gives an isomorphism
$L_h:H^n(X)\to H^{n+2}(X)$; its inverse $\beta$ is algebraic because
the Lefschetz standard conjecture holds for abelian varieties. Let $e$
be the algebraic primitive projector in degree $n$ on $X$, with image
$\ker L_h^2$, and extend it by zero in other degrees. Define
\[
 r=\beta j_*,\qquad \Pi_F=j^*e\beta j_*.
\]
Since $j_*j^*=L_h$, one has $rj^*=1$ and $\Pi_F^2=\Pi_F$ on
cohomology. The correspondence $\Pi_F$ is algebraic. The incidence
detection lemma from JAC-T4.1 therefore induces an idempotent on
$P^m(Y)$ in the isogeny category.

The decomposition
\[
 H^n(Y)=j^*H^n(X)\oplus\ker j_*
\]
is orthogonal for the middle pairing on $Y$. On its first summand the
pairing is $\int_X uvh$. For $u\in\ker L_h^2$ and
$v=L_hw\in L_hH^{n-2}(X)$ one has
$\int_Xuvh=\int_X(L_h^2u)w=0$. Thus $\Pi_F$ is self-adjoint for
the pairing on $Y$. Saito's transpose compatibility and nondegenerate
incidence duality imply that its image $B_F$ is a nondegenerate
self-dual isogeny factor of $P^m(Y)$. Moreover it is primitive for $Y$:
if $u\in\ker L_h^2$, then
$j_*L_{j^*h}j^*u=L_h^2u=0$, and
$j_*:H^{n+2}(Y)\to H^{n+4}(X)$ is injective by dual weak Lefschetz.

It remains to verify the sign on all curves in $B_F$. Pull a normalized
rational Poincar\'e family back to a smooth pointed curve $C$ mapping to
$B_F$, and apply the parameter degree-one projector. Its class is an
algebraic tensor
\[
 \xi\in H^1(C)\otimes F.
\]
Set $\eta=(1\otimes r)\xi\in H^1(C)\otimes H^n(X)$.
It is algebraic, $j^*\eta=\xi$, and $h^2\eta=0$.
If $b$ is the genus of $C$ and $\theta$ the canonical principal
polarization on $J(C)$, the algebraic inverse to the Abel pullback
lifts $\eta$ to an algebraic class
$w\in H^1(J(C))\otimes H^n(X)$. It satisfies $h^2w=0$.
For $H=\theta+h$ on the abelian variety $J(C)\times X$,
\[
 H^{b+1}w=0.
\]
Indeed terms with $h$-exponent at least two vanish, while terms with
$h$-exponent zero or one have excessive degree on $J(C)$.
Consequently $w$ is primitive of codimension $m$. In its quadratic form
only one term of $H^b$ survives, and the Abel curve-class formula yields
\[
 \begin{aligned}
 \int_{J(C)\times X}w^2H^b
 &=b\int_{J(C)\times X}w^2\theta^{b-1}h\\
 &=b!\int_{C\times X}\eta^2h
 =b!\int_{C\times Y}\xi^2.
 \end{aligned}                                                \tag{H9.9}
\]
The factor $b!$ is positive. By 032, Corollary 8.4, the last intersection
has sign $(-1)^m$ and is zero only if $\xi=0$. The same curve test and
nef-plus-nondegenerate criterion as in JAC-T4.3 therefore prove that
the signed incidence class on $B_F$ is ample. This finishes the proof.
\hfill$\square$

The projectors used here implement the hyperplane formula of cote 16,
scan 40; the cohomological fixed/variable decomposition belongs to
SGA 7 II, XVIII, \S5. The input 032 is applied only to the genuinely
abelian product $J(C)\times X$, never to $J(C)\times Y$.

\subsection*{H9-T5. A finite-field Newton witness eliminates the remaining part}

\textbf{Claim.} Let a Lefschetz pencil of ample sections of an abelian
$2m$-fold be defined over a finite field of characteristic $p$. Write
$U$ for its smooth locus and $\mathcal V_\ell$ for its variable
$R^{2m-1}$-local system. Assume that $\mathcal V_\ell$ is nonzero and
remains absolutely irreducible under every finite-index restriction of
geometric monodromy, for a fixed prime $\ell\ne p$. Let
$B_{\rm var}$ be the variable incidence factor of $P^m(Y)$ over the
generic field, after a finite extension making the factor and universal
cycle defined. If one smooth closed fibre has a Frobenius eigenvalue on
its variable cohomology with normalized $p$-adic slope outside
$[m-1,m]$, then $B_{\rm var}=0$. In that case H9-T4 proves the primitive
Theorem 4 for the full incidence quotient of the generic fibre.

\textbf{Proof.}
The algebraic correspondence $1-j^*L_h^{-1}j_*$ induces the variable
idempotent on the incidence quotient. It is self-adjoint, so its image
$B_{\rm var}$ has nondegenerate incidence duality. A normalized Poincar\'e
cycle gives, for every $\ell\ne p$, an algebraic correspondence map
\[
 V_\ell B_{\rm var}\longrightarrow\mathcal V_\ell(m)
                                                        \tag{H9.10}
\]
over a finite cover of $U$. This map is injective: a vector in its
kernel pairs to zero with every other vector under the pullback of the
middle pairing, whereas that pullback is the realization of the
nondegenerate incidence isogeny. A choice of polarization on
$B_{\rm var}$ identifies $V_\ell B_{\rm var}$ with
$H^1(B_{\rm var},\mathbf Q_\ell)(1)$. Thus (H9.10) is equivalently
\[
 H^1(B_{\rm var},\mathbf Q_\ell)(-(m-1))
       \lhook\joinrel\longrightarrow\mathcal V_\ell.       \tag{H9.11}
\]
If $B_{\rm var}$ is nonzero, the irreducibility hypothesis, which survives
the finite field-of-definition extension, forces (H9.11) to be an
isomorphism.

At every point of the finite cover lying above a smooth fibre, the
right side is unramified. Therefore the left side is unramified and
the N\'eron--Ogg--Shafarevich criterion gives good reduction of
$B_{\rm var}$ there. Smooth proper base change specializes (H9.11).
If the residue field has cardinality $q$, the Frobenius eigenvalues on
$H^1$ of this reduction are algebraic integers $\gamma$ for which
$q/\gamma$ is also an algebraic integer. Hence, at every embedding into
$\overline{\mathbf Q}_p$,
\[
 0\leq\frac{v_p(\gamma)}{v_p(q)}\leq1.
\]
The Tate twist $-(m-1)$ multiplies geometric-Frobenius eigenvalues by
$q^{m-1}$. Thus all slopes on the variable part must belong to
$[m-1,m]$. Passing to a finite extension of a residue field raises
each eigenvalue and $q$ to the same positive power and leaves these
normalized slopes unchanged. A single smooth original fibre with a
slope outside the interval therefore contradicts $B_{\rm var}\ne0$.
This proves the claim.\hfill$\square$

The geometric irreducibility hypothesis is the standard monodromy
condition studied in SGA 7 II, XVIII, 6.7.2; it is stated explicitly
here so that no exceptional pencil is silently included. The good
reduction criterion and the integrality/duality of Frobenius on abelian
varieties are background inputs. This proof does not assert that every
high-degree pencil has the required Newton witness. Producing that
witness is the residual concrete gate of this method; purity of the
complex absolute values alone is insufficient to do so.


% Draft extension submitted for independent root/programme review.
% It becomes a checked finding only when the full argument is read and logged.

\subsection*{H9-T6. The variable incidence sign from integral monodromy}

The Newton test in H9-T5 is a useful way to prove that the variable
incidence factor is zero. The following argument instead determines its
sign when it is nonzero. It uses Grothendieck's integral monodromy pairing,
not a sign assigned to one $\ell$-adic scalar.

\textbf{Hypotheses.} Let $Y$ be the generic fibre of an odd-dimensional
Lefschetz pencil of fibre dimension $n=2m-1$. Suppose its middle
cohomology has an algebraic self-adjoint variable projector, and let
$B$ be the corresponding incidence factor of $P^m(Y)$. Assume:
\begin{enumerate}
\item for every $\ell\ne\operatorname{char}k$, the variable local system
      $\mathcal V_\ell$ is nonzero and absolutely irreducible after every
      finite-index restriction of geometric monodromy;
\item at one ordinary quadratic singular fibre the vanishing cycle in
      $\mathcal V_\ell$ is nonzero for every such $\ell$;
\item $B$ has its nondegenerate symmetric incidence duality and a
      normalized rational Poincar\'e cycle, as in Saito's construction.
\end{enumerate}
For the ample-section pencils considered in H9-T4, the algebraic variable
projector is $1-j^*L_h^{-1}j_*$ on the relevant odd degree. The first two
hypotheses are the precise Picard--Lefschetz monodromy hypotheses; their
standard generic-pencil instances are in SGA 7 II, XVIII, 6.7.2.

\textbf{Claim.} Under these hypotheses the rational divisor class
\[
 \ell_m=\frac{(-1)^m}{2}\Delta_B^*d_B
\]
is a rational polarization of $B$, in every characteristic. Consequently,
under the same monodromy hypotheses, H9-T4 and H9-T6 together prove the
primitive Theorem 4 for the full incidence quotient of a generic ample
section of an abelian variety, accepting 032 Corollary 8.4 for its fixed
part.

\paragraph{H9-L1: a rational square detected at almost all primes.}
If $r\in\mathbf Q^*$ is a square in $\mathbf Q_\ell$ for all but
finitely many primes $\ell$, then $r$ is a square in $\mathbf Q$.
In particular $r>0$.

\emph{Proof.} Otherwise the quadratic field $\mathbf Q(\sqrt r)$
defines a nonprincipal quadratic Dirichlet character $\chi$ of some
conductor. At primes outside the conductor, local squareness implies
$\chi(\ell)=1$, so only finitely many Euler factors of $L(s,\chi)$
differ from those of $\zeta(s)$, for real $s>1$. The quotient
$L(s,\chi)/\zeta(s)$ therefore has a finite positive nonzero limit as
$s\downarrow1$. On the other hand, a nonprincipal periodic character
has bounded partial sums and zero sum over a period. Summation by parts
shows that $\sum_{n\ge1}\chi(n)n^{-s}$ stays bounded as $s\downarrow1$
(indeed converges at $s=1$), whereas $\zeta(s)\to+\infty$.
This is a contradiction. No Riemann hypothesis is used.\hfill$\square$

\paragraph{H9-L2: the rank-one monodromy coefficient.}
Let an abelian variety have semistable reduction over a strictly
henselian trait, and suppose the toric rank is one. Let $L_0$ be a
genuine polarization and let $E_0$ be the alternating form on the Tate
module defined by $c_1(L_0)$. Choose a primitive generator $v_0$ of the
toric cocharacter lattice, and choose the same generator of
$\mathbf Q_\ell(1)$ in the monodromy operator and in the values of $E_0$.
Then the tame monodromy logarithm has the form
\[
 N_\ell(x)=-q_0 E_0(x,v_0)v_0,\qquad q_0\in\mathbf Q_{>0},
                                                        \tag{H9.12}
\]
where the rational number $q_0$ is independent of $\ell$.

\emph{Derivation from the cited theorem.} SGA 7 I, IX, 9.2 identifies
the monodromy operator with the scalar extension of the integral
pairing between the two toric character lattices. Theorem 10.4(a),
printed p.~444, asserts that this integral pairing is common to all
primes; 10.4(b) asserts positivity after composing with a genuine
polarization. In rank one, choose integral bases and write $u$ for
the monodromy coefficient and $s\ne0$ for the polarization coefficient
on the two lattices. Positivity is $us>0$. In compatible rational toric
and quotient coordinates, $N(y)=u v_0$ and $E_0(v_0,y)=s$, so
\[
 q_0=u/s=(us)/s^2>0.
\]
This proves rationality and positivity simultaneously; it never interprets
one $\ell$-adic square as a real square. The sign in (H9.12) is the minus
sign from alternating the arguments of $E_0$. It agrees with the curve
case: a one-node Jacobian has $N(y)=v_0$ with
$E_0(v_0,y)=1$. The compatibility of this curve normalization with
the general monodromy pairing is explicitly SGA 7 I, IX, 12.5 and 12.7.
Under a ramified extension of traits of index $e$, the monodromy pairing
is multiplied by $e>0$ (IX, 10.3.5, printed p.~443).

\paragraph{Proof of H9-T6.}
If $B=0$ there is nothing to prove, so suppose $B\ne0$. Work over one
finite extension of the generic field on which $B$, the normalized
Poincar\'e cycle and a genuine polarization $L_0$ are defined. Let
\[
 f_\ell:V_\ell B\longrightarrow\mathcal V_\ell(m)
\]
be the map of realizations given by this cycle. As in H9-T5,
nondegenerate incidence duality makes $f_\ell$ injective. The
irreducibility hypothesis makes it an isomorphism, for every
$\ell\ne\operatorname{char}k$.

Every geometric endomorphism of $B$ is defined over a finite extension.
Tate-module faithfulness gives
\[
 \operatorname{End}^0(B_{\overline K})\otimes\mathbf Q_\ell
   \lhook\joinrel\longrightarrow
 \operatorname{End}_{G'}(V_\ell B)
\]
after a sufficiently small finite-index monodromy subgroup $G'$.
Absolute irreducibility makes the right side $\mathbf Q_\ell$.
It follows that $\operatorname{End}^0(B_{\overline K})=\mathbf Q$.
Consequently the nondegenerate symmetric class $\ell_m$ is
\[
 \ell_m=c L_0\quad\text{in }\operatorname{NS}(B)_{\mathbf Q},
 \qquad c\in\mathbf Q^* .                              \tag{H9.13}
\]
The task is to determine the sign of this single rational number.

Let $Q_\ell$ denote the middle cup-product pairing on
$\mathcal V_\ell(m)$, with values in $\mathbf Q_\ell(1)$.
The incidence formula gives the following exact sign:
\[
 E_{\ell_m}=(-1)^{m-1}f_\ell^*Q_\ell=cE_0.            \tag{H9.14}
\]
To verify it, write the mixed class of the universal cycle as
$z=\sum_i e_i\otimes u_i$, with $e_i\in H^1(B)$ and
$u_i\in H^n(Y)(m)$. In the product $z_{13}z_{23}$ the odd-degree
class $u_i$ crosses the degree-one parameter class $e_j$, producing
a minus sign. Thus
\[
 d_B=-\sum_{i,j}Q_\ell(u_i,u_j)e_i\boxtimes e_j.
\]
Pulling back to the diagonal and dividing by two shows that the
alternating form of $\Delta^*d_B/2$ is $-f_\ell^*Q_\ell$.
Multiplying by $(-1)^m$ proves (H9.14).

At the chosen singular fibre, take a finite trait extension that makes
the abelian variety semistable and the toric lattice split. Let its
ramification index relative to the original pencil parameter be $e>0$.
The nonzero Picard--Lefschetz operator has rank one. Hence the toric rank
of $B$ is one: by IX, 9.2 and the nondegeneracy in 10.4, rank of monodromy
equals rank of the toric character lattice.

SGA 7 II, XV, 3.3.3--3.3.6, printed pp.~191--192, supplies a vanishing generator
$\delta_\ell$, normalized up to sign by the primitive quotient of the
exceptional quadric. With the same tame-character convention as in
(H9.12), its formula, in relative dimension $2m-1$, is
\[
 N_\ell(x)=e(-1)^m Q_\ell(x,\delta_\ell)\delta_\ell.
                                                        \tag{H9.15}
\]
Here twists are suppressed only after choosing a generator of
$\mathbf Q_\ell(1)$. More intrinsically, XV defines the untwisted
vanishing generator in $H^{2m-1}(Y)(m-1)$; tensoring it with that
generator puts it in $H^{2m-1}(Y)(m)$, and dividing both cup-product
values and the tame character by that same generator gives (H9.15).
Thus no unrecorded unit is introduced into the coefficient.

Put $v_\ell=f_\ell^{-1}\delta_\ell$. Since $f_\ell$ intertwines
monodromy, $v_\ell$ spans the toric line, so
$v_\ell=a_\ell v_0$ for some $a_\ell\in\mathbf Q_\ell^*$.
Combining (H9.14) and (H9.15) gives
\[
 N_\ell(x)=-e c a_\ell^2 E_0(x,v_0)v_0.
\]
Comparison with (H9.12), using $E_0$ nondegenerate, proves
\[
 a_\ell^2=\frac{q_0}{e c}\quad\text{for every }\ell\ne
 \operatorname{char}k.                                \tag{H9.16}
\]
The right side is the same nonzero rational number for all these primes.
H9-L1 therefore says it is a rational square and in particular positive.
Since $q_0>0$ and $e>0$, one obtains $c>0$. Equation (H9.13) now proves
that $\ell_m$ is a rational polarization.

For the final assertion, the fixed primitive and variable incidence
summands are orthogonal and their dualities are block diagonal.
H9-T4 polarizes the former and the argument just completed polarizes
the latter. A product of the two rational polarizations is ample,
and ampleness descends under the isogeny giving the decomposition.
This proves the claimed extension of Theorem 4.\hfill$\square$

\paragraph{Scope.}
The result concerns the geometric generic fibre under explicit monodromy
hypotheses. It does not assert positivity on every special fibre of a
pencil, where new incidence factors can appear. Nor does it identify
the height of a merely homologically trivial generic cycle with a
N\'eron--Tate height on $P^m$: the E^4 domain obstruction remains.
The original H9(b) therefore does not follow from this extension.
\part{Arithmetic proof drafts and unresolved gates}
This part contains full arithmetic proof drafts, including conditional
lemmas whose hypotheses are substantive. The current checking status is
recorded by finding ID in the event log. Inclusion here does not assert
that the remaining arithmetic height-sign problem has been solved.
\section{Arithmetic sibling of cote 44: corrected theorem and remaining
gate}\label{arithmetic-sibling-of-cote-44-corrected-theorem-and-remaining-gate}

Run: \texttt{JAC1964-20261121}. Author/checker:
\texttt{/root/arithmetic\_sign}, GPT-6 Astra Pro; effort setting not
exposed. Date: 2026-11-21 UTC. These are working proof sections for
incorporation into the parent manuscript. They are not an independent
review of the parent manuscript. No Lean build was attempted here.
Accepted attached results are used at their stated scope.

\subsection{A. Conventions and exact
background}\label{a.-conventions-and-exact-background}

Let \(K\) be a number field, \(A/K\) an abelian variety of dimension
\(g\), and \(\ell=c_1(L)\) for a symmetric ample line bundle. All Chow
groups and Chow motives have rational coefficients. Write \(L\) also for
multiplication by \(\ell\). The raw Beilinson--Bloch pairing, normalized
by \([K:\mathbf Q]\), is denoted \(h_K\). Thus \(h_K\) is invariant
under finite extension of the field of definition. This normalization is
distinct from the choice whether the bilinear Néron--Tate form is the
polarization of \(\widehat h\) with or without a factor two; below the
Néron--Tate convention is Künnemann's.

\textbf{Background AR-B1 (existence and functoriality).} Künnemann,
\emph{Height pairings for algebraic cycles on abelian varieties}, Ann.
Sci. ENS 34 (2001), Theorem 1.6 and Corollary 1.7, give the pairing on
every abelian variety over a number field. Section 1.3 records
compatibility with correspondences; §1.4, equation (7), gives the
extension-of-field scaling before normalization. The pairing uses
extensions perpendicular to vertical cycles; an arbitrary closure on a
model is insufficient. The correspondence projection formula is also
Künnemann 1996, §5, equation (18), from Lemma 4.1.

\textbf{Background AR-B2 (Picard comparison and sign).} Künnemann,
\emph{Higher Picard varieties and the height pairing}, Amer. J. Math.
118 (1996), Theorem 8.2, equation (23), identifies the height of
algebraically trivial cycles with the Néron--Tate biextension pairing on
higher Picard varieties, with the positive comparison integer \(k_X^p\).
His Theorem 6.1 calibrates the raw curve height as \textbf{minus} the
positive Jacobian pairing. Let \(B^p(A_K)\) be algebraically trivial
cycles modulo geometric incidence equivalence. Theorem 12.1(ii) gives

\[
 x\ne0,\quad L^{g+2-2p}x=0
 \quad\Longrightarrow\quad
 (-1)^p h_K(x,L^{g+1-2p}x)>0.
\]

The old regular-model hypothesis in the 1996 statement can be met after
finite extension using AR-B1; the positive base-change factor does not
affect the conclusion. This is a known theorem, not a new consequence of
GRH or of the attached CM Hodge theorem. For \(g=2m-1,p=m\), the
positive candidate is

\[
q_A(x,y)=(-1)^m h_K(x,y),\qquad Lx=Ly=0.
\tag{AR.1}
\]

Consequently ``the BB height equals a positive Néron--Tate height''
requires a specified signed duality. With the parent's convention
\(\phi_L(a)=t_a^*L\otimes L^{-1}\), the signed incidence is
\((-1)^m D\). Its Hodge cup-product counterpart carries \((-1)^{m-1}\);
the incidence dictionary supplies the extra minus. The \(m=1\)
Faltings--Hriljac case is an indispensable check.

\textbf{Background AR-B3 (motivic Lefschetz).} Künnemann, \emph{A
Lefschetz decomposition for Chow motives of abelian schemes}, Invent.
Math. 113 (1993), Theorem 5.1, Theorem 5.2, and Corollary 5.3, all on
printed p.96, supply actual Chow projectors, motivic hard Lefschetz, and
transposition compatibility. For \(g=2m-1\), the primitive projector
\(\pi_P\) on \(h^g(A)\) satisfies \({}^t\pi_P=\pi_P\); the primitive
motive \(P\) has its perfect middle Poincaré pairing. These three
statements were checked directly in the rendered source page, not
inferred from a cohomological decomposition.

\textbf{Background AR-B4 (degree one is faithful).} The functor from
abelian varieties up to isogeny to degree-one Chow motives is fully
faithful (contravariantly). In particular,

\[
\operatorname{End}_{\mathrm{Chow}}(h^1B)=\operatorname{End}(B)_{\mathbf Q}^{\mathrm{op}},
\]

and the Betti realization of this endomorphism algebra is faithful. A
directly inspected modern statement is Laga--Shnidman, \emph{Ceresa
cycles of bielliptic Picard curves}, arXiv:2312.12965, Proposition 4.5,
citing Kings, \emph{Higher regulators, Hilbert modular surfaces, and
special values of L-functions}, Duke Math. J. 92 (1998), Proposition
2.2.1. The contravariant/opposite convention changes no invertibility
assertion below.

\subsection{B. Exact removal of the irrelevant Beauville
grades}\label{b.-exact-removal-of-the-irrelevant-beauville-grades}

For \([n]^*x=n^{2p-s}x\), say \(x\) has Beauville grade \(s\). The
well-known height selection rule is

\[
h_K(\operatorname{CH}^p_{(s)},\operatorname{CH}^{g+1-p}_{(t)})=0
\quad\text{unless}\quad s+t=2.
\tag{AR.2}
\]

An explicit literature locator is Zhang, \emph{Gross--Schoen cycles and
dualising sheaves}, Invent. Math. 179 (2010), §5.3, printed p.68. For
completeness its proof is one line: the projection formula gives

\[
h_K([n]^*x,[n]^*y)=n^{2g}h_K(x,y),
\]

whereas the eigenvalue computation gives \(n^{2(g+1)-s-t}\) times the
original pairing. Take \(n>1\).

\textbf{Proposition AR-1 (primitive grade reduction).} Put

\[
V^p=\ker\bigl(L^{g+2-2p}:\operatorname{CH}^p(A)_{\mathrm{hom}}
\longrightarrow\operatorname{CH}^{g+2-p}(A)\bigr),\qquad 2p\le g+1.
\]

Then \(V^p_{(s)}=0\) for \(s\ge2\). Every \(V^p_{(s)}\) with \(s\le0\)
belongs to the radical of the signed height on \(V^p\). Thus
semipositivity on the full primitive group is equivalent to
semipositivity on grade one. Strict positivity on the full group is
equivalent to strict positivity on grade one together with the vanishing
of the primitive grades \(s\le0\).

\emph{Proof.} A grade-\(s\) class of codimension \(p\) is a Chow class
of \(h^i(A)\) with \(i=2p-s\). For \(s\ge2\), put \(e=g+2-2p\). Then
\(i\le g\) and \(e\le g-i\). Motivic hard Lefschetz makes \(L^{g-i}\) an
isomorphism. If \(L^e x=0\), multiply by \(L^{g-i-e}=L^{s-2}\) to obtain
\(L^{g-i}x=0\), hence \(x=0\). Consequently every surviving primitive
grade is at most one. Multiplication by a symmetric \(\ell\) has grade
zero. Equation (AR.2) therefore pairs a primitive grade \(s\le0\) only
with grade \(2-s\ge2\), which has just been proved zero. This proves the
radical assertion and both equivalences. No assertion that all negative
Beauville grades vanish was used. \(\square\)

For \(g=2m-1,p=m\), the grade-one component of \(V^m\) is exactly

\[
V^m_{(1)}=\operatorname{CH}^m(P).
\tag{AR.3}
\]

Indeed the motivic Lefschetz decomposition of \(h^g(A)\) is a direct sum
of the primitive middle summand and longer Lefschetz strings. On the
middle-degree term of each longer string, \(L\) has a left inverse; on
\(P\), it is zero. This checks that (AR.3) is a Chow-level assertion,
rather than merely a statement about realization.

\textbf{Corollary AR-1a (the threefold has no additional lower-grade
gate).} If \(g=3,p=2\), then \(\operatorname{CH}^2_{(s)}(A)=0\) for
every \(s<0\), and \(\operatorname{CH}^2_{(0)}(A)_{\mathrm{hom}}=0\).
Consequently the full primitive homologically trivial group equals
\(\operatorname{CH}^2(P)\), and strict positivity has no separate
lower-grade vanishing condition.

\emph{Proof.} Fourier gives \[
\mathcal F:\operatorname{CH}^2_{(s)}(A)
 \simeq\operatorname{CH}^{1+s}_{(s)}(\widehat A).
\] For \(s<-1\) the target has negative codimension. For \(s=-1\) it is
\(\operatorname{CH}^0_{(-1)}=0\). For \(s=0\), a homologically trivial
class maps to a homologically trivial divisor of grade zero. Such a
divisor lies in \(\operatorname{Pic}^0\), which has grade one, so it is
zero. Fourier is invertible and preserves homological triviality. AR-1
removes all primitive grades at least two; only grade one remains.
\(\square\)

\subsection{C. A CM Chow summand with a strictly positive
height}\label{c.-a-cm-chow-summand-with-a-strictly-positive-height}

Assume the attached 032 Theorem 1.1, i.e.~the rational Hodge conjecture
for products of CM abelian varieties. Its Lemma 8.1 supplies effective
CM realizations; in weight one it is the usual realization by a CM
abelian variety. The statement below uses those accepted inputs,
classical motivic Lefschetz, and the established Picard height theorem.
It does not use GRH.

\textbf{Theorem AR-2 (orthogonal Picard summand and exact remaining
arithmetic gate).} Let \(A/K\) be CM, \(\dim A=2m-1\), with symmetric
ample \(L\). After a finite number-field extension \(K'/K\), there are a
CM abelian variety \(B/K'\), an actual self-transpose Chow idempotent
\(e\) of the primitive middle motive \(P\), and an orthogonal
decomposition

\[
P=M\oplus N,\qquad M\simeq h^1(B)(1-m),
\tag{AR.4}
\]

with the following properties.

\begin{enumerate}
\def\labelenumi{\arabic{enumi}.}
\tightlist
\item
  \(H(M)\) is the maximal rational sub-Hodge structure of \(H(P)\)
  having only types \((m,m-1)\) and \((m-1,m)\). Thus \(H(N)\) has no
  nonzero such sub-Hodge structure.
\item
  \(\operatorname{CH}^m(M)\simeq\operatorname{Pic}^0(B)(K')_{\mathbf Q}\).
  The restriction of \(q_A=(-1)^m h_{K'}\) to this full Chow summand is
  a Néron--Tate form for a rational polarization, with the positive
  comparison normalization of AR-B2. It is strictly positive on every
  nonzero class.
\item
  The two Chow summands are height orthogonal. Therefore the
  semipositivity problem on the full primitive homologically trivial
  Chow group over \(K'\) reduces \textbf{exactly} to
\end{enumerate}

\[
(-1)^m h_{K'}(z,z)\ge0
\qquad(z\in\operatorname{CH}^m(N)).
\tag{AR-GATE}
\]

For strict positivity one must additionally prove the vanishing of the
primitive grades \(s\le0\) and strict positivity on the nonzero classes
of \(\operatorname{CH}^m(N)\).

\emph{Proof.} Let \(W\subset H(P)\) be the sum of all rational sub-Hodge
structures of the indicated types. Polarizable Hodge structures are
semisimple; this sum is a sub-Hodge structure and has an orthogonal
complement. Because \(A\) is CM, so is \(W\). Its twist \(W(m-1)\) is
effective of weight one. The degree-one case of 032 Lemma 8.1, followed
by Poincaré complete reducibility, gives a CM \(B\) and an isomorphism

\[
H^1(B)(1-m)\simeq W.
\]

A complex CM abelian variety admits a number-field model (Milne, Abelian
Varieties, version 1.10, Proposition 16.10, printed p.65). After
enlarging \(K\), both \(B\) and chosen polarizations are defined over
the same number field. Write \(b=\dim B\). The inclusion of this
weight-one twist into \(H^g(A)\) is a Hodge class of codimension
\(b+m-1\) on \(B\times A\). Theorem 1.1 applied to \(B\times A\) makes
it an algebraic correspondence. This is exactly the codimension
bookkeeping in the proof of 032 Corollary 8.2: \(n=1,r_0=m-1\).

The correspondence can be chosen over \(\overline{\mathbf Q}\). Indeed
spread each of its finitely many cycle components over a finitely
generated extension of \(\overline{\mathbf Q}\), or equivalently
represent it by a complex point of the Hilbert scheme of the fixed
variety \(B\times A\). Take the connected component of that finite-type
Hilbert scheme containing this point. Because the scheme is defined over
the algebraically closed field \(\overline{\mathbf Q}\), that component
contains a \(\overline{\mathbf Q}\)-point. The universal family is flat,
and its Betti cycle class on the fixed complex variety is locally
constant on the analytification; connectedness therefore preserves the
class. Specializing each component in this way preserves the original
rational Hodge correspondence. Finitely many algebraic coefficients and
cycle equations require only one finite extension \(K'\). Sandwich the
resulting correspondence between the actual degree-one and primitive
Chow projectors. We obtain a Chow morphism

\[
u:M_0=h^1(B)(1-m)\longrightarrow P
\]

whose realization is the chosen injection onto \(W\).

There is an algebraic, perfect alternating form on \(M_0\) given by a
polarization of \(B\) and motivic hard Lefschetz. On \(P\), use
\((-1)^{m-1}\) times the middle Poincaré form. Its restriction to the
types occurring in \(W\) is a Hodge polarization. Both forms are perfect
at the Chow-motive level: for \(B\) this is motivic hard Lefschetz; for
\(P\) it is AR-B3. They define a Chow adjoint \(u^\dagger\). The
realization of

\[
a=u^\dagger u\in\operatorname{End}_{\mathrm{Chow}}(M_0)
\]

is positive and invertible: for the positive Hodge Hermitian forms,
\((av,v)=(uv,uv)>0\) for nonzero \(v\). By AR-B4 this is an actual
rational isogeny endomorphism of \(B\), so \(a^{-1}\) exists in the Chow
endomorphism algebra. This step does \textbf{not} infer Chow
invertibility from a general cohomological isomorphism; it uses full
faithfulness specifically in degree one.

Set

\[
v=a^{-1}u^\dagger,\qquad e=uv=u a^{-1}u^\dagger.
\]

Then \(vu=1\), \(e^2=e\), and \(e^\dagger=e\), by direct composition.
Since the form on \(P\) is a scalar multiple of its Poincaré pairing,
\(e^\dagger=e\) says precisely \({}^t e=e\) within \(P\). Its
realization is the orthogonal projector onto \(W\). The complementary
motive \(N=(P,1-e)\) therefore has the stated realization. No assertion
that a realization-zero complementary motive is zero is made; such a
conclusion would need a separate nilpotence theorem. This avoids a
hidden lift from homological to Chow motives.

The retraction identifies \(\operatorname{CH}^m(M)\) with
\(\operatorname{CH}^1(h^1B)=\operatorname{Pic}^0(B)(K')_{\mathbf Q}\).
Such classes are algebraically trivial. This identification is also
injective modulo incidence equivalence: if \(u_*x\) is incidence
equivalent to zero, apply the correspondence \(v\); then \(x\) is
incidence equivalent to zero as a divisor on \(B\), hence \(x=0\) in
\(\operatorname{Pic}^0(B)_{\mathbf Q}\). The latter assertion follows
either from the identity correspondence in the definition of incidence
equivalence, or from the ordinary Picard universal property. Thus
AR-B2's primitive strict positivity applies to every nonzero class of
this Chow summand, not only to an unspecified quotient. Its Theorem 8.2
identifies the form with the pullback of the signed primitive Picard
polarization, giving the stated Néron--Tate description.

For \(x=e_*x\) and \(y=(1-e)_*y\), the height projection formula gives

\[
h_{K'}(x,y)=h_{K'}(x,{}^t e(1-e)_*y)=0.
\]

Combining this with AR-1 and (AR.3) proves the exact gate and the
strict-positivity assertion. Finite base change is harmless for classes
descending to \(K\): normalized heights are unchanged, and pullback on
rational Chow groups is injective since pushforward composed with
pullback multiplies by \([K':K]\). \(\square\)

\textbf{Scope caution.} The decomposition is asserted over \(K'\). For a
fixed initial \(K\), positivity on all of \(N(K')\) is a sufficient
condition for positivity on descended classes, not a claim of logical
equivalence with positivity over \(K\) alone. The exact equivalence in
the theorem concerns the group over \(K'\). The theorem is a derived
consequence of stated inputs; no claim of literature novelty is made.

\subsection{D. A small explicit complement which the reduction does not
erase}\label{d.-a-small-explicit-complement-which-the-reduction-does-not-erase}

Let \(E/K\) have an automorphism \(\rho\) of order six acting on
\(H^{1,0}(E)\) by a primitive sixth root \(\zeta\), and assume this
automorphism is defined over \(K\). For example take a \(j=0\) elliptic
curve after extending its number field. On \(A=E^3\), use the product
polarization and let \(\rho_\Delta=(\rho,\rho,\rho)\). Let
\(\Pi=\pi_1\otimes\pi_1\otimes\pi_1\) be the degree-one tensor Chow
projector. Put

\[
e_3=\Pi\,\frac16\sum_{r=0}^5(-1)^r\rho_\Delta^{*r}.
\tag{AR.5}
\]

\textbf{Proposition AR-3.} Equation (AR.5) is an actual rational Chow
idempotent, self-transpose on the middle motive, and its motive \(T\) is
primitive. Its Hodge realization has dimension two and types
\((3,0),(0,3)\). Its entire \(\operatorname{CH}^2(T)\) lies in Beauville
grade one and is homologically trivial. Consequently its height is a
concrete surviving case of (AR-GATE), with positive-sign candidate
\(h_K(z,z)\).

\emph{Proof.} The cyclic group character projector is idempotent because
the coefficient character \(r\mapsto(-1)^r\) is a rational character of
the order-six group. It commutes with \(\Pi\). Transposition replaces
group elements by their inverses and fixes the character projector;
\(\pi_1\) of an elliptic curve is self-transpose. In the tensor
\(H^1(E)^{\otimes3}\), the eigenvalues are \(\zeta^3=-1\), \(\zeta\),
\(\zeta^{-1}\), and \(\zeta^{-3}=-1\), according to the number of
holomorphic factors. The two all-holomorphic/all-antiholomorphic tensors
are exactly the \(-1\)-eigenspace. Multiplication by each factor's
divisor class kills the factor \(h^1(E)\), since it would land in
\(h^3(E)=0\); this is exact in the Chow-motive grading. Hence
\(L\Pi=0\), so \(T\) is primitive. Finally \(T\subset h^3(A)\), and
\(\operatorname{CH}^2(h^3(A))=\operatorname{CH}^2_{(1)}(A)\). It has
zero ordinary degree-four cycle realization. \(\square\)

This is a specified motive and a specified height problem. It is not an
assertion that \(\operatorname{CH}^2(T)\) is known to contain a nonzero
class, nor an assertion that all its heights have been determined. No
general height-to-\(L'\) identity for this group follows from the
supplied papers.

\subsubsection{D.2 The complementary gate contains a nonzero
Fermat-quartic
cycle}\label{d.2-the-complementary-gate-contains-a-nonzero-fermat-quartic-cycle}

\textbf{Theorem AR-7 (a nonzero CM complementary class).} Let
\(C:x^4+y^4=z^4\), let \(J=\operatorname{Jac}(C)\), and choose a
degree-one base class after a finite number-field extension. Write \[
\kappa=[C]-[-1]_*[C]\in\operatorname{CH}^2(J)_{\rm hom,\mathbf Q}.
\] Let \(\pi_P\) be the primitive middle Chow projector for a symmetric
theta divisor, and let \(e\) be the Picard-summand projector of AR-2.
Then \[
z_N=(\pi_P-e)_*\kappa\in\operatorname{CH}^2(N)
\tag{AR.9}
\] is nonzero modulo algebraic equivalence, even over
\(\overline{\mathbf Q}\). In particular (AR-GATE) is a genuine remaining
sign problem on a group containing a specified nonzero class. In this
example \(H(N)\) has precisely types \((3,0),(0,3)\), each of dimension
one.

\emph{Proof.} Bloch, \emph{Algebraic cycles and values of L-functions},
J. reine angew. Math. 350 (1984), Lemma (3.1), p.99, gives an isogeny
\(J\to E^2\times E'\), of degree a power of two, where \(E:y^2=x^3-x\)
and \(E':y^2=x^3+x\). Thus \(J\) is CM; over \(\overline{\mathbf Q}\) it
is isogenous to the cube of a \(j=1728\) elliptic curve. Bloch's Theorem
(4.1), p.104, says that this Ceresa cycle has infinite order modulo
algebraic equivalence. Its class modulo algebraic equivalence is
independent of the chosen embedding (Theorem (3.4) and following text,
p.102).

We first check that \(\kappa\) has grade one, so the Chow projector
\(\pi_3\) does not discard its nonzero Griffiths class. It is odd under
\([-1]^*\), hence only odd Beauville grades can occur. The Fourier
isomorphism \[
\mathcal F:\operatorname{CH}^2_{(s)}(J)
 \simeq\operatorname{CH}^{1+s}_{(s)}(\widehat J)
\] kills every odd \(s\ge3\), by dimension, and every odd \(s\le-3\), by
negative codimension. For \(s=-1\) the target is
\(\operatorname{CH}^0_{(-1)}(\widehat J)=0\), since
\(\operatorname{CH}^0\) is spanned by the fundamental class and has
grade zero. Therefore
\(\kappa\in\operatorname{CH}^2_{(1)}(J)=\operatorname{CH}^2(h^3J)\).

Motivic Lefschetz gives \(h^3(J)=P\oplus Lh^1(J)\). The codimension-two
Chow group of the second summand is an image under \(L\) of
\(\operatorname{Pic}^0(J)_{\mathbf Q}\); all these classes are
algebraically trivial. The same is true of \(\operatorname{CH}^2(M)\),
by AR-2. Consequently \[
\kappa-z_N=(1-\pi_P)_*\kappa+e_*\kappa
\] is algebraically trivial. Bloch's theorem now forces \(z_N\) to be
nonzero modulo algebraic equivalence. This argument is unchanged after
finite extension.

For the final Hodge-type assertion, the rational CM decomposition of
\(H^3(E^3)\) splits off its two extreme types \((3,0),(0,3)\). One can
see rationality explicitly by applying the scalar CM endomorphism
\(1+i\): its eigenvalues on the extreme types satisfy \(T^2+4T+8\),
while those on \((2,1),(1,2)\) satisfy \(T^2-4T+8\). These coprime
rational polynomials give the two rational Hodge substructures. The full
level-one part has dimension \(18\). The Lefschetz summand \(LH^1(J)\)
is six-dimensional and lies entirely in that part. The extreme
two-dimensional substructure is already primitive, by Hodge type. Thus
the primitive middle piece has a twelve-dimensional maximal level-one
part and the indicated two-dimensional complement. \(\square\)

This proof refutes a CM-wide assertion that homological triviality is
algebraic triviality up to torsion: the rational Griffiths quotient is
nonzero already for this CM Jacobian. It does not refute the arithmetic
Hodge-index sign conjecture. The new statement here is the precise
placement of the known nonzero cycle in the orthogonal complementary
Chow summand, using AR-2; Bloch's nonvanishing theorem is an accepted
established input.

The Beauville Fourier grading used above is the standard Fourier
decomposition (Beauville, \emph{Sur l'anneau de Chow d'une variété
abélienne}, Math. Ann. 273 (1986), 647--651). An explicitly read
statement with the grading convention is Shen--Vial, \emph{The Fourier
transform for certain hyperkähler fourfolds}, Introduction A, p.2,
displayed (1) and property (a). Its Fourier inversion is also recalled
in Künnemann 1993, §2, printed p.88.

\subsubsection{D.3 A canonical nonzero complementary cycle with an exact
local height
formula}\label{d.3-a-canonical-nonzero-complementary-cycle-with-an-exact-local-height-formula}

\textbf{Theorem AR-8 (Fermat symmetry and Gross--Schoen reduction).} For
the Fermat quartic choose the hyperflex \(P_0=(1:0:1)\) and its Abel
embedding \(f(x)=[x-P_0]\). Let \(X_1\) be the grade-one part of
\(f_*[C]\), and set \[
\eta=\mathcal F(X_1)\in\operatorname{CH}^2(J).
\] One can choose the splitting in AR-2 so that
\(\eta\in\operatorname{CH}^2(N)\). This class is nonzero modulo
algebraic equivalence. Let \(K'\) contain the relevant definitions and
admit a semistable model of \(C\), and put \[
\Omega_{K'}=\frac{\omega_a^2}{[K':\mathbf Q]},
\qquad
\Phi_{K'}=\frac{1}{[K':\mathbf Q]}
 \sum_v\varphi(C_v)\deg(v).
\] Then the raw normalized middle height has the exact value \[
h_{K'}(\eta,\eta)
 =\frac16 h_{K'}(\Delta_{P_0},\Delta_{P_0})
 =\frac7{24}\Omega_{K'}-\frac16\Phi_{K'}.
\tag{AR.10}
\] Thus this specified nonzero complementary class has nonnegative
height exactly when \(7\Omega_{K'}\ge4\Phi_{K'}\), and strictly positive
height exactly when the inequality is strict. There is no positive
Picard contribution to subtract for this choice of splitting.

\emph{Proof.} The tangent \(x=z\) meets \(C\) at \(P_0\) with
multiplicity four, so \(K_C\sim4P_0\). Let \(G=\mu_4^2\) act by the two
coordinate scalings on the affine equation \(x^4+y^4=1\). Each
\(\gamma P_0\) satisfies \(4\gamma P_0\sim K_C\), so \(\gamma P_0-P_0\)
is four-torsion in \(J\). Translation by a torsion point acts
identically on rational Chow groups: if \(na=0\), then
\([n]\circ t_a=[n]\); since \([n]_*\) is invertible on rational Chow
groups by the Chow--Künneth eigenvalue decomposition,
\([n]_*t_{a*}=[n]_*\) implies \(t_{a*}=1\). Therefore the class
\(f_*[C]\), its grade-one part, and its Fourier transform \(\eta\) are
\(G\)-invariant. The last assertion uses the \(G\)-invariance of the
symmetric principal theta class; the Fourier kernel is consequently
\(G\)-equivariant.

The holomorphic differentials have \(G\)-characters \[
(1,1),\quad(1,2),\quad(2,1)\qquad\text{in }(\mathbf Z/4)^2.
\] Their exterior triple has trivial character. In a mixed exterior
triple of type \((2,1)\), canceling a holomorphic and antiholomorphic
character leaves one of the three nonzero characters above; using three
distinct indices gives \((0,2),(2,0),(2,2)\), again nonzero. Conjugation
handles type \((1,2)\). Hence \(H^3(J)^G\) consists exactly of the two
extreme Hodge types. In particular it has no overlap with \(W=H(M)\).

Let \(p_G=|G|^{-1}\sum_{\gamma\in G}\gamma_*\). This is an exact
self-transpose Chow idempotent. It commutes with the primitive projector
because \(G\) preserves theta. In AR-2 replace the constructed inclusion
by \[
u'=(1-p_G)u.
\] Its realization is unchanged on \(W\); therefore the same adjoint
argument gives a split inclusion and a self-transpose projector \(e'\).
Now \(p_Gu'=0\) holds exactly, and transposition gives
\((u')^\dagger p_G=0\), hence \(e'p_G=0\). Since \(p_G\eta=\eta\), we
get \(e'\eta=0\). Zhang's Theorem 1.5.6 says \(\eta\) is primitive;
therefore \(\eta\in\operatorname{CH}^2(N)\) for this splitting.

By AR-7, the Ceresa class is \(2X_1\) and is nonzero modulo algebraic
equivalence. Fourier is an invertible algebraic correspondence and
preserves algebraic equivalence, so \(\eta\) is likewise nonzero. Zhang,
\emph{Gross--Schoen cycles and dualising sheaves} (2010), Theorem 1.5.6,
printed p.14, gives the first equality in (AR.10) in genus three. His
Theorem 1.3.1, printed p.5, gives the second because \(P_0-K_C/4=0\) in
the rational Picard group. Dividing every global intersection by
\([K':\mathbf Q]\) gives the displayed normalization. \(\square\)

The theorem makes the remaining calculation concrete; it does not
establish the inequality in (AR.10). Good finite places contribute zero
to \(\Phi_{K'}\). For the displayed integral plane quartic, every odd
prime is a smooth-reduction prime by the partial derivatives. Thus only
places over \(2\) and the archimedean places can contribute to
\(\Phi_{K'}\). The admissible self-intersection \(\Omega_{K'}\) also
includes arithmetic metric and model information. A computation of these
terms, including the semistable reduction at \(2\), would settle the
height of this cycle. Positivity of this one nonzero class would still
not establish positivity of the whole group \(\operatorname{CH}^2(N)\).

\subsection{E. Mixed-characteristic and constant-family
checks}\label{e.-mixed-characteristic-and-constant-family-checks}

The correct arithmetic analogue of an intersection formula has the form

\[
h_K(z,z')=\frac1{[K:\mathbf Q]}\widehat{\deg}
 \bigl(\widehat z^{\mathrm{adm}}\cdot\widehat z'^{\mathrm{adm}}\bigr),
\tag{AR.6}
\]

with vertical-orthogonal extensions and the compatible Green-current
data. Künnemann 2001, §1.2--1.4, is the exact construction used here.
Thus a raw closure intersection does not become a height merely by
inserting \((-1)^m\). The parent's quadric counterexample to the
uncorrected pencil formula is consistent with this distinction: integral
fibers do not remove all higher-codimension local correction terms.

\textbf{Constant-family check.} For the constant family
\(B\times C\to C\) over a field \(k\), and a constant section \(s_P\)
with \(P\in B(k)\), the geometric canonical height over \(k(C)\) is
zero. Indeed the pullback of a line bundle from \(B\) along the constant
map is a constant one-dimensional vector space tensored with
\(\mathcal O_C\), hence has degree zero; the same holds after every
multiplication by \(n\). If \(k=K\) is a number field and \(P\) is
nontorsion, its arithmetic Néron--Tate height on \(B/K\) is strictly
positive for an ample polarization. Therefore the arithmetic value
cannot be recovered from that geometric degree alone. The missing term
is genuine arithmetic metric/finite-place information. This is a
counterexample to any claimed direct replacement of a geometric height
by the number-field height with all these terms omitted, not a
counterexample to a correctly normalized biextension formula.

For ordinary divisors, Faltings--Hriljac gives \(h_K(D,D)<0\) when the
Jacobian class of \(D\) is nontorsion. For a primitive codimension-two
class on a threefold the proposed positive sign is \(+h_K\). These two
low-dimensional checks prevent an unsigned ``BB is NT-positive''
statement.

\subsection{F. What GRH actually
gives}\label{f.-what-grh-actually-gives}

\textbf{Lemma AR-4 (central sign).} Let \(F\) be holomorphic in a
neighborhood of the real interval \([c,\infty)\), real on that interval,
positive for sufficiently large real arguments, and nonzero for every
real \(s>c\). If \(F\) vanishes to order \(r\ge1\) at \(c\), then
\(F^{(r)}(c)>0\). In particular \(F'(c)\ge0\), with strict inequality
exactly when \(r=1\).

\emph{Proof.} Continuity and the absence of a real zero on
\((c,\infty)\) give \(F(s)>0\) there. In the expansion
\(F(c+t)=a_rt^r+O(t^{r+1})\), \(a_r\) is real and nonzero. Positivity
for small \(t>0\) forces \(a_r>0\), and \(F^{(r)}(c)=r!a_r\).
\(\square\)

If \(F=G L\), with \(G(c)>0\) real and \(L(c)=0\), then
\(F'(c)=G(c)L'(c)\). Thus GRH supplies the lemma's zero-free interval
for an appropriately normalized real self-dual completed \(L\)-function,
provided poles and the positivity normalization have also been handled.
A complex non-self-dual Hecke \(L\)-function has no ordered derivative
sign until an appropriate real combination has been specified.

\textbf{Corollary AR-5 (conditional height sign).} If a proven formula
gives \(q_A(z,z)=C L'(c)\), \(C>0\), with the hypotheses of AR-4, then
\(q_A(z,z)\ge0\). It is strictly positive if the central zero is simple.
If the central zero has higher order, the formula gives zero; it does
not prove \(z=0\). To deduce positivity for an entire cycle space one
needs a spanning family and control of cross terms, or a formula for
every vector in that space. A regulator determinant alone is
insufficient: the real symmetric matrix
\(\operatorname{diag}(-1,-1,1,\ldots,1)\) has positive determinant and
is not positive definite.

The attached Alpöge \texttt{Statements.lean} quantifies over Dirichlet
characters \(\chi:\operatorname{DirichletCharacter}(\mathbf C,N)\). Its
\texttt{Main.lean} proves exactly those advertised Dirichlet statements
from imported results; it contains no theorem for arbitrary number-field
Hecke characters. We accept the cited theorem; this is a scope check,
not a re-verification of its proof. The Masullo index labels 14.1.2.1
(printed p.529) a \textbf{Conjecture}, and 14.1.2.3 (p.531) a
proposition conditional on that conjecture and 14.1.2.2. The row on
p.522 whose opening reads ``13.3.1.2 (Divisor of L(s, chi))'' includes
unfinished statement text (the index parses its number as 13.3.1).
Accordingly any use of broad Hecke GRH is explicitly \textbf{conditional
(Hecke-GRH/Masullo)}. No such input supplies the height identity in
AR-5.

\subsubsection{F.2 The concrete CM derivative and the newer
triple-product
comparison}\label{f.2-the-concrete-cm-derivative-and-the-newer-triple-product-comparison}

Bloch 1984, Proposition (5.2), printed p.107, proves that the Hecke
\(L\)-function of the Fermat-quartic rank-two motive of types
\((3,0),(0,3)\) has a simple zero at \(s=2\). This is a theorem about
order of vanishing. The proposition does not state an arithmetic height
identity or an ordered sign for the derivative. The introduction,
printed p.95, explicitly presents the height formula as a possible
further assertion. Thus the supplied Dirichlet-GRH theorem, Bloch's
nonvanishing, and (AR.10) are three distinct inputs.

A stronger next comparison is available in Yuan--Zhang--Zhang,
\emph{Triple product L-series and Gross--Kudla--Schoen cycles},
manuscript dated 14 November 2023, with appendix by Yifeng Liu.
Proposition 1.3.2, printed p.7, expresses the trilinear cycle height as
a coefficient \(\mathcal L(\pi)\) times a product of local invariant
pairings. Theorem 1.3.3, printed p.8, states \[
\mathcal L(\pi)=
\frac{16\zeta_F(2)^2}{L(1,\pi_L,\mathrm{ad})}
 L'(1/2,\pi_L)+\sum_{p\in S}c_p\log p,
\qquad c_p\in\overline{\mathbf Q},
\tag{AR.11}
\] where \(S\) consists of rational primes below ramification places of
\(\pi\). The exact generalized Gross--Kudla formula holds in their
everywhere-unramified case. This is stronger than the older September
2012 manuscript, Theorem 1.2.4, whose hypotheses included at least two
finite nonsplit quaternion places.

The Fermat quartic is \(X_0(64)\); for example the explicit modular-form
embedding and the relation \(x^4+y^4=z^4\) are given in Fang-Ting Tu,
\emph{Arithmetic of Quaternion Orders and Its Applications} (Waseda
dissertation, February 2012), §5.3.2, p.59. Hence the relevant modular
elliptic constituents are unramified away from \(2\). If one carries out
the precise trilinear-cycle identification and the local test-vector
normalizations for this example, the possible correction in (AR.11) is
supported at \(2\). Its shape is \(c_2\log2\); its vanishing or sign is
not supplied by an \(L\)-function zero-free theorem. The whole
identification is retained as a task, rather than treating a general
modular parametrization as an automatic equality with \(\eta\).

The newer Gao--Zhang manuscript, \emph{Heights of Ceresa and
Gross--Schoen cycles}, Theorems 1.1 and 1.4, constructs canonical adelic
height objects and proves a lower bound on a suitable Zariski-open part
of moduli. This offers genuine geometry for the Bost route below. A
bound with an additive constant on an unspecified open subset does not
determine the sign at this fixed Fermat point. The specific formula
(AR.10), or a completed comparison including the term in (AR.11),
remains the concrete height gate.

\subsection{G. Testing the two suggested Alpöge
mechanisms}\label{g.-testing-the-two-suggested-alpuxf6ge-mechanisms}

\subsubsection{G.1 Sign-only Weil/explicit-formula
route}\label{g.1-sign-only-weilexplicit-formula-route}

The note's footnote 3, printed p.3, explains the sign in the
finite-field trace calculation using concentration in middle cohomology
and the opposing zero-mode parity, with a bound from repeated Parseval.
Step 5 (pp.3--4) changes to a von Mangoldt/Guinand--Weil variance setup.
Neither statement identifies the cycle intersection or arithmetic height
form with the analytic form.

A minimal algebra check isolates the missing information. Let \(F=pI\)
on \(\mathbf Q^2\), \(p>0\). Both \(b_+=I\) and
\(b_\pm=\operatorname{diag}(1,-1)\) satisfy \(F^tbF=p^2b\); they have
the same Frobenius eigenvalues and traces and different signatures. Thus
eigenvalue purity and trace parity do not determine a compatible form's
signature. This is an information-loss test, not a proposed geometric
counterexample. A valid geometric transfer still needs an identification
of the form and its polarization.

\subsubsection{G.2 Bost slope route}\label{g.2-bost-slope-route}

Bost, \emph{Algebraic leaves of algebraic foliations over number
fields}, Publ. Math. IHES 93 (2001), §4.1.4, Proposition 4.5, equation
(4.11), gives for an injective \(K\)-linear map between hermitian
bundles

\[
\widehat{\deg}_{n}(E)
\le\operatorname{rk}(E)\bigl(\widehat\mu_{\max}(F)+h(\varphi)\bigr).
\tag{AR.7}
\]

Proposition 4.6, equation (4.14), is its filtered version. The note's
step 13, p.4, names this inequality. The attached Lean map points to
\texttt{Formalism/LatticeSlopes/LatticeSlopesProofs/SlopeInequality.lean}
and \texttt{Formalism/LatticeSlopes/Statement.lean}; these
implementation files were \textbf{not attached or read}, and no claim is
made to have run them.

\textbf{Proposition AR-6 (precise slope gate).} Suppose that a geometric
construction associates to a cycle \(z\) a hermitian line
\(\overline{\mathcal B}_z\) with
\(\widehat{\deg}_n\overline{\mathcal B}_z=q_A(z,z)\), and an injective
map \(\varphi:\mathcal B_z^\vee\to F_z\) satisfying

\[
\widehat\mu_{\max}(F_z)+h(\varphi)\le0.
\tag{AR.8}
\]

Then \(q_A(z,z)\ge0\); a strict upper bound in (AR.8) gives strict
positivity.

\emph{Proof.} Apply (AR.7) to the rank-one source
\(\overline{\mathcal B}_z^\vee\), whose normalized degree is
\(-q_A(z,z)\). \(\square\)

The hypotheses in AR-6 are substantive: the line must be the cycle's
canonical arithmetic biextension/height object, and the map and norm
bounds must be proved. Merely declaring a trivial line to have metric
\(e^{-q_A(z,z)}\) is circular. To see why the inequality by itself
yields no sign, let \(E_a=(\mathcal O_K,\|1\|_\sigma=e^{-a})\). Then
\(\widehat{\deg}_n(E_a)=\widehat\mu_{\max}(E_a)=a\), and the identity
map has height zero. Equation (AR.7) becomes \(a\le a\) for every
positive or negative real \(a\).

This explicitly tests the next tool: a Bost proof of (AR-GATE) would
require (AR.8) for the height line of every class in the complementary
motive \(N\), including local and archimedean terms. Neither the
four-page note nor the supplied Lean map provides this bridge. On the
Picard summand the Poincaré biextension already gives the canonical
object and the result is known independently.

\subsection{H. Sources and reading
record}\label{h.-sources-and-reading-record}

\begin{itemize}
\tightlist
\item
  Attached 032: Theorem 1.1; Lemma 8.1; Corollary 8.2 and its proof.
  Accepted inputs; targeted statement/codimension reading.
\item
  Attached cote 44 scan004: Theorem 5 and its margin. It assumes
  algebraically trivial generic restrictions; its printed intersection
  formula is handled separately by the parent's H9 counterexample and
  correction.
\item
  Künnemann 1993:
  https://gdz.sub.uni-goettingen.de/download/pdf/PPN356556735\_0113/LOG\_0011.pdf
  . Direct PDF downloaded, rendered printed pp.96--98; Theorems 5.1,5.2,
  Corollary 5.3 directly read.
\item
  Künnemann 1996: DOI https://doi.org/10.1353/ajm.1996.0033 . Primary
  paper's indexed PDF at
  https://citeseerx.ist.psu.edu/document?doi=24752c170be79846ab861e9c80a8977cac60035f\&repid=rep1\&type=pdf
  . Web retrieval exposed full extracted text after direct-download
  failures. Lemma 4.1, equation (18), Theorems 6.1,8.2,9.1,12.1, and
  Remark 12.2 were targeted-read; not a full proof audit.
\item
  Künnemann 2001:
  https://www.numdam.org/item/10.1016/s0012-9593(01)01068-0.pdf .
  §1.2--1.4, Theorem 1.6, Corollary 1.7 directly read via extracted
  primary PDF.
\item
  Milne, Abelian Varieties v1.10:
  https://www.jmilne.org/math/CourseNotes/AV110.pdf . Proposition 16.10,
  printed p.65, model over a number field; statement read, proof omitted
  in this source.
\item
  Zhang 2010:
  https://web.math.princeton.edu/\textasciitilde shouwu/publications/gross-schoen.pdf
  . §5.3, p.68, height selection; directly read.
\item
  Laga--Shnidman: https://arxiv.org/pdf/2312.12965 . Proposition 4.5,
  Proposition 5.9 and Proposition 5.10 directly read; Kings Prop.2.2.1
  is cited through this source, not directly read.
\item
  Bost 2001:
  https://www.numdam.org/item/PMIHES\_2001\_\_93\_\_161\_0.pdf . §4.1.4,
  Propositions 4.5--4.6, equations (4.11),(4.14), directly read.
\item
  Attached Alpöge note: pp.3--4, footnote3, steps5 and13; attached
  Main/Statements; Lean-map §4 and §5. No formal build and no claim to
  inspect omitted implementation files.
\item
  Attached Masullo index: rows 13.3.1.2,14.1.2.1,14.1.2.3. The full
  paper is not attached; broad Hecke GRH is retained as an explicit
  hypothesis.
\end{itemize}

\subsection{I. Self-check summary}\label{i.-self-check-summary}

AR-1: exponent \(g-i-e=s-2\) checked; grade zero of symmetric \(\ell\)
checked; no unproved Beauville vanishing inserted. AR-2: codimension
\(b+m-1\), twist \(1-m\), Chow degree \(m+(1-m)=1\), exact inverse
specifically in \(h^1\), self-transpose projector, incidence injection,
field extension scope, and height orthogonality each checked. AR-3:
idempotent group algebra character, transpose/inverse, eigenvalue
multiplicities \(1,3,3,1\), primitive tensor action, and grade
\(2p-i=1\) checked. AR-4: poles, real-valuedness, positivity
normalization, order of zero, and derivative/leading-coefficient
distinction checked. AR-6: dual-source sign and rank-one rescaling test
checked. AR-7: Fourier odd-grade elimination, primitive versus Lefschetz
summand, algebraic triviality of the discarded terms, and the two CM
characteristic polynomials checked; root independently confirmed this
argument. AR-1a: all negative grades and the homologically trivial
grade-zero divisor image checked. AR-8: hyperflex canonical class,
torsion translation using invertible multiplication, all twenty exterior
character triples, exact averaged-projector annihilation, and the
constants 7/24 and 1/6 checked. The character enumeration independently
returned only the all-holomorphic and all-antiholomorphic triples as
invariant.

Remaining parent-facing review request: independently check AR-2's
motivic-adjoint sign and field-descent argument before promoting it in
the final paper. It is a complete written argument from stated
background, not a formal certificate and not a claimed solution of
(AR-GATE).
\section{What remains to be proved}
The full primitive arithmetic height form on the complementary CM Chow
motive remains to be controlled. A formula for one special cycle does
not prove positivity for every cycle. Analytic derivative signs require
an applicable real self-dual $L$-function and a proved height identity
with its local terms and positive constants. The Bost slope inequality
requires an actual canonical height line and a proved norm estimate.
Neither missing input is silently supplied by purity or Dirichlet GRH.

The geometric H9 assertions cannot be completed in their original form:
the counterexamples above refute the formula and the asserted incidence
domain. Their repaired statements are the explicitly stated admissible
formula, no-correction criterion, and conditional quotient lemma.
\begin{thebibliography}{99}
\bibitem{Saito1979} H. Saito, \emph{Abelian varieties attached to cycles of intermediate dimension}, Nagoya Math. J. 75 (1979), 95--119. Definition 3.1; Proposition 3.4; Theorem 3.6; \S4.1; Proposition 4.4; Corollary 4.5; Theorem 4.6; Proposition 5.2. \url{https://doi.org/10.1017/S0027763000024855}.
\bibitem{Kleiman68} S. L. Kleiman, \emph{Algebraic cycles and the Weil conjectures}, in Dix expos\'es sur la cohomologie des sch\'emas (1968), 359--386, Theorem 2A11.
\bibitem{MilneAV} J. S. Milne, \emph{Abelian Varieties}, course notes, v2.00 (2008), Chapter I, Theorem 11.1; Chapter III, Theorem 6.8 and Summary 6.11. \url{https://www.jmilne.org/math/CourseNotes/AV.pdf}.
\bibitem{MumfordAV} D. Mumford, \emph{Abelian Varieties}, Oxford University Press, 1970. Chapter II, \S6, Application 1, p.~60; Chapter III, \S16, Index Theorem, p.~155.
\bibitem{ConradAV} B. Conrad, \emph{Abelian varieties}, notes by T. Feng (2015), Example 4.2.14, p.~38.
\bibitem{SGA7I} A. Grothendieck et al., \emph{Groupes de monodromie en g\'eom\'etrie alg\'ebrique, I}, SGA 7 I, Lecture Notes in Mathematics 288, Springer, 1972. Expos\'e IX, \S9.2, 10.3.5 (p.~443), Theorem 10.4 (p.~444), \S\S12.5 and 12.7.
\bibitem{SGA7II} P. Deligne and N. Katz, eds., \emph{Groupes de monodromie en g\'eom\'etrie alg\'ebrique, II}, SGA 7 II, Lecture Notes in Mathematics 340, Springer, 1973. Expos\'e XV, 3.3.3--3.3.6, pp.~191--192; XVII, \S2.2; XVIII, 5.3.2--5.3.5, pp.~287--288, Proposition 6.2, p.~319, Corollary 6.4, 6.7.2; XX, Theorem 4.1 and Remark 4.3, pp.~355--356.
\bibitem{Bloch1984L} S. Bloch, \emph{Algebraic cycles and values of L-functions}, J. reine angew. Math. 350 (1984), 94--108. Lemma 3.1, p.~99; Theorem 4.1, p.~104; Proposition 5.2, p.~107.
\end{thebibliography}
\end{document}
